% darwin-memo: arXiv preprint skeleton.
% Compiles with: pdflatex main && bibtex main && pdflatex main && pdflatex main
% Neutral preprint style (article); swap to a venue style (neurips/icml/colm)
% by replacing the class + geometry block only. No exotic packages.
% These are local review layouts. Neither implements a selected venue template
% or establishes compliance with a submission page limit.
\ifdefined\bodyonly
  \documentclass[10pt,twocolumn]{article}
\else
  \documentclass[11pt]{article}
\fi

\usepackage[utf8]{inputenc}
\usepackage[T1]{fontenc}
\usepackage{geometry}
\ifdefined\bodyonly
  \geometry{letterpaper, margin=0.75in, columnsep=0.25in}
\else
  \geometry{a4paper, margin=1in}
\fi
\usepackage{amsmath, amssymb}
\usepackage{booktabs}
\usepackage{array}
\usepackage{multirow}
\usepackage{algorithm}
\usepackage{algpseudocode}
\usepackage{xcolor}
\usepackage{graphicx}
\usepackage{subcaption}
\usepackage{pgfplots}
\pgfplotsset{compat=1.18}
\usepackage[numbers,sort&compress]{natbib}
\usepackage[hidelinks]{hyperref}
\usepackage{microtype}
\usepackage{enumitem}

% This paper is float-heavy: nine tables and two figures in the body, all
% full-width. LaTeX's defaults reserve so much of each page for text that
% double-column floats queue up and spill onto pages of their own, which
% costs pages a venue limit counts. These are the standard relaxations for
% a table-dense paper; nothing is scaled or shrunk.
\renewcommand{\topfraction}{0.92}
\renewcommand{\dbltopfraction}{0.92}
\renewcommand{\bottomfraction}{0.7}
\renewcommand{\textfraction}{0.06}
\renewcommand{\floatpagefraction}{0.8}
\renewcommand{\dblfloatpagefraction}{0.8}
\setcounter{topnumber}{4}
\setcounter{dbltopnumber}{4}
\setcounter{totalnumber}{6}

\newcommand{\todo}[1]{\textcolor{red}{[\textbf{TODO:} #1]}}
\newcommand{\resscale}{\sigma}      % resource scale
\newcommand{\credit}{g}             % credit gain
\newcommand{\upkeep}{u}             % upkeep
\newcommand{\Emax}{E_{\max}}        % energy cap

% One source, two builds. Both carry every word; they differ only in where
% the supporting detail sits.
%
%   tectonic main.tex                        -> preprint: detail inline
%   tectonic body.tex (\def\bodyonly{})      -> venue: detail in the appendix
%
% A block named \inplace{x} lives in sections/detail/x.tex and is typeset
% either where it is written or in the appendix's deferred region, never in
% both and never in neither. tests/test_paper_deferred_blocks.py holds that
% invariant: a block dropped from one build is the failure this mechanism
% would otherwise make silent.
%
% The venue page limit is measured to where the appendix begins, not to the
% last page -- security venues exclude the appendix -- so both builds are
% complete documents and neither deletes anything.
% \pointer{...} is the sentence that only makes sense once the detail has
% moved: it names where the block went, and vanishes in the build where the
% block is on the same page.
\ifdefined\bodyonly
  \newcommand{\inplace}[1]{}
  \newcommand{\pointer}[1]{#1}
  % Two columns are narrower than one, so the tables step down a size.
  \newcommand{\tabsize}{\footnotesize}
\else
  \newcommand{\inplace}[1]{\input{sections/detail/#1}}
  \newcommand{\pointer}[1]{}
  \newcommand{\tabsize}{\small}
\fi

\title{\textbf{Attacking the Curator:\\
Curation-Targeted Attacks on Agent Memory,\\
and What Survives Them}}

\author{
  Roger Sim\~oes\\
  \texttt{roger@wearegoodai.com}\\
  Good AI
  % \and add co-authors here
}

\date{\today}

\begin{document}
\maketitle

% The metadata abstract is maintained in abstract-arxiv.txt.
\begin{abstract}
We study feedback corruption against agent-memory curators: an attacker
changes the outcome reports that determine which lessons are retained, while
an untouched ground-truth record measures harm and useful-memory loss. Our
contributions are an explicit budgeted threat model, a controlled comparison of
curation policies and their failure boundaries, and a reusable artifact with
committed observations and offline table reconstruction.

In a filesystem-backed simulation with a modeled restoration penalty, a bounded
outcome-credit ledger preserves benign capability at one corrupted report per
cycle and nearly preserves it at two. At four reports per cycle, its mean benign
capability falls to $0.06$. Simple counters fail earlier in the initial grid,
but a wider sweep finds forgiveness settings that retain useful memory under
small attack budgets. We retain this counterexample to universal superiority.
These results measure behavior under the stated simulator and feedback model;
they do not measure physical restoration cost or establish causal lesson value.

Historical coding-agent experiments provide a separate, less conclusive view:
no memory arm establishes an advantage over carrying no memory, and attack
responses vary by repository. They lack the complete billing and prospective
noninferiority evidence needed to claim lower total operating cost with preserved
task success. We distinguish offline replay from runs where memory changes
subsequent agent behavior, document missing evidence, and provide separate paths
for inexpensive reconstruction and experiment reruns.
\end{abstract}

\section{Introduction}
\label{sec:intro}

A frozen language model can be given new knowledge without retraining by
attaching an external memory: a store of past experience retrieved into
the context window at inference time \citep{packer2023memgpt,
park2023generative, shinn2023reflexion, zhao2024expel}, or in its strong
form a dedicated memory \emph{model} while the executive LLM stays fixed
\citep{quek2026memo}. Every such system faces the same question once it
runs long enough: \emph{what gets to stay?} Memory that only grows is
eventually dominated by stale, redundant or actively wrong entries, and
the failure is not hypothetical---memory banks an LLM continuously
rewrites degrade over time, with useful entries becoming faulty as the
rewriting accumulates \citep{zhang2026faulty}.

The prevailing answer puts a \emph{model in the loop} to curate: an LLM
rates importance \citep{park2023generative}, rewrites trajectories
\citep{zhang2026faulty}, or judges what to keep. These share two costs.
They are expensive---every curation decision is a model call---and they
introduce a \emph{proxy} between the agent and the world. A judge grades
plausibility, not consequences, and a proxy that can be optimized against
can be gamed: the reward-hacking and semantic-drift failure mode
self-training systems fall into when no external criterion anchors
quality \citep{dodgson2026survival}.

\paragraph{The curator is an attack surface, and nobody has measured it.}
Whatever plays the curator---a classifier, a judge, a heuristic, a
ledger---it decides on some signal, and that signal is reachable by an
adversary. The memory-security literature defends the store at
\emph{write} (does this text look malicious?) and at \emph{retrieval}
(should this entry be injected?) \citep{mpbench2026, memsecbench2026},
and both read content. Neither asks what happens when the input to the
\emph{deciding rule} is corrupted. The consequence is not that poison
survives; it is worse. A defence that removes entries on bad evidence,
fed bad evidence, removes the defender's own working memory. The
adversary need inject nothing at all, and the more decisively a
mechanism curates the more thoroughly it destroys the store it protects.
We call this a \emph{curation-targeted attack}.

The threat has been named. \citet{lin2026selfevolving} describe an
adversary who corrupts the performance signal to control ``the long-term
composition of the agent's cognitive state without ever writing to the
memory store directly,'' and \citet{lin2026memsurvey} give forgetting
its own lifecycle phase. Both also record that it has not been
measured---``relatively underexplored'' in the first, ``architecturally
plausible but empirically unstudied'' in the second
\citep{lin2026memsurveyv1}, whose later version drops that assessment
rather than revising it. That is the gap we
close, and it is worth being plain that closing it is the contribution:
this paper does not discover that curators are attackable, it makes the
threat quantitative.

This paper operationalizes it, measures six curation mechanisms under it,
and reports which survive. That requires a curator whose deciding signal
is not a model's opinion. Our starting observation is that a large and
practically important class of agent tasks comes with a \emph{conserved,
physically measurable resource} that already settles value without anyone
grading anything: a storage-management agent frees or wastes real bytes,
a coding agent makes a test suite pass or fail, an issue-resolution agent
resolves a GitHub issue or does not. In each case the environment returns
a scalar that is \emph{conserved}---you cannot manufacture freed bytes or
passing tests by being persuasive---and so cannot be hacked by a more
eloquent memory. The resource is conserved; the channel that
\emph{reports} it is ordinary software, and \S\ref{sec:whocanlie} says
who can write to that channel without controlling the resource. Building
on the environment-mediated selection principle of
\citet{dodgson2026survival}, we make that scalar the only currency of
memory survival.

\paragraph{\textsc{darwin-memo}.}
Each memory entry carries an \emph{energy} balance. Every cycle, every
entry pays a fixed upkeep, so existence is not free. An entry earns energy
back only when a task it decided (or supported) produces a positive
measured resource delta, with credit flowing along provenance and bounded
by a $\tanh$ so no single outcome makes an entry immortal and no single
disaster instantly kills one that has been right many times. When an
entry's energy crosses zero it is buried. The consequences fall out of the
dynamics rather than from any rule about quality: trivia nobody consults
starves on upkeep; a poisoned entry that advises a harmful action is
\emph{executed by the damage it causes} as the negative delta flows back
to it; near-duplicate survivors consolidate, pooling energy. There is no
reward model, no LLM judge, and no human curation anywhere, and because
settlement is a scalar the environment already returns, it is essentially
free.

\paragraph{The result, and its boundary.}
Under attack the mechanisms separate cleanly (\S\ref{sec:adversary}), and
what makes the difference is a bounded credit buffer with earn-back: a
lie costs energy an entry can win back, where a counter spends a life it
cannot. At one corrupted measurement per cycle the ledger keeps all of
its benign capability where every strike and quarantine policy we test
keeps at most about half, and a Thompson-sampling bandit keeps its
capability only by never removing anything, so retention is not defence.
The ledger is not exempt either, and we publish its boundary rather than
stopping at the win.

\paragraph{What we claim, and what we do not.}
It would be easy, and wrong, to claim that conserved-resource selection
dominates all alternatives. It does not, and the failures are not
marginal: a one-line heuristic matches the ledger on outcomes and, with a
real model deciding adoption, revokes poison five to nine cycles
\emph{sooner} (\S\ref{sec:wef}); a strike counter and a quarantine
policy each beat it on either side of a noise threshold in a second
environment family; a Thompson-sampling retrieval bandit reconstructed
from \citet{xu2026ael} matches it under one-sided noise; and on real
software-engineering tasks it shows no learning-curve advantage over a
token-matched random control (\S\ref{sec:swebench}). We report all of
them, and the honest summary is narrow: \emph{outside an adversarial
regime, the ledger's extra machinery buys leanness and little else.} Our
second contribution is therefore the \emph{regime map}, a pre-registered
account of which curator to reach for when.

\paragraph{Contributions.}
\begin{enumerate}[leftmargin=1.2em, itemsep=1pt, topsep=2pt]
\item \textbf{An operational threat model for an attack the literature
has named but not measured.} In a \emph{curation-targeted attack} the
adversary injects no poison and instead corrupts the signal a curator
decides on, turning the defence into the instrument that starves the
defender's benign memory. We give it a resource-bounded adversary, an
accounting identity and a budget axis, so it becomes an experiment rather
than a scenario (\S\ref{sec:adversary}), and \S\ref{sec:whocanlie} says
who can reach the settlement channel without controlling the resource
behind it.
\item \textbf{The first measurement of how curation mechanisms fare under
it}---six mechanisms at five budgets over 30 seeds---establishing that a
bounded credit buffer with earn-back is what separates the survivors,
that retention without removal is not a defence, and that every
mechanism including ours has a boundary, which we locate and publish.
\item \textbf{A label-free, judge-free curation mechanism} that grounds
survival in a conserved environmental resource, with credit assignment
along provenance and a formal description (\S\ref{sec:method}).
\item \textbf{A measured retention attack on usage-based curation},
closing a second gap the literature names without measuring
\citep{lin2026memsurvey}. A \emph{query-only} adversary---one who never
writes, never settles, and crosses no channel our threat model
names---drives two of the three components of a usage-based retention
score by asking questions, and thereby owns roughly the last $13\%$ of
the store's lifetime ($16$ of $16$ cells; $0$ of $32$ under the flat
upkeep we ship). A counterfactual swapping one denominator isolates the
cause: the exploitable property is not that the signal is \emph{usage}
but that it is \emph{relative} (\S\ref{sec:potentiation}).
\item \textbf{A sweep of the baseline family's own free parameter, and
the refutation it produced} (\S\ref{sec:countersweep}). A baseline nobody
tuned is not a baseline.
\item \textbf{The regime map} (\S\ref{sec:regime}), with control arms
reconstructed from the literature---a retrieval bandit
\citep{xu2026ael}, an LLM-judge settler, and a Generative-Agents-style
salience score \citep{park2023generative}---run on the same stdlib-only
harness rather than argued about, under pre-registered hypotheses,
seed-level bootstrap intervals, exact paired permutation tests, and
committed per-seed data validated against a manifest in CI.
\end{enumerate}

\section{Related Work}
\label{sec:related}

\paragraph{Memory poisoning, and where its defences sit.}
\citet{mpbench2026} systematize memory-poisoning attacks on LLM agents
and report the split that motivates this work: every off-the-shelf
detector they test loses between $18$ and $42$ points of true-positive
rate moving from strong-signal payloads to weak-signal ones---payloads
written to read like ordinary operational facts, carrying, in their
reading, ``no syntactic anomaly''. \citet{memsecbench2026} track poisoning through a lifecycle
and identify selective repair as the open stage: removing poison without
destroying benign memory alongside it. We adopt their
Write--Execute--Forget checkpoint structure as a protocol
(\S\ref{sec:wef}), reimplemented rather than ported, and make no
comparison to their numbers.

What both share with the wider defence literature is \emph{where} the
defence sits: at write, screening text before it is stored, or at
retrieval, screening entries before they are injected. Both read content,
and content is exactly what a weak-signal payload is constructed to
survive. Our threat model is orthogonal to that axis rather than a better
position on it: we do not ask whether the defence can recognize the
payload, but what it does when the evidence it decides on has itself been
tampered with.

\paragraph{Attacking the curator: named, and left unmeasured.}
We are not the first to describe this threat, and we are explicit about
that because the distinction between naming it and measuring it is the
whole of our contribution here. \citet{lin2026selfevolving} catalogue
the attack surface of self-evolving agent systems and state the model
directly: an adversary who can ``shift embedding proximity or corrupt
the performance signal gains indirect control over the long-term
composition of the agent's cognitive state without ever writing to the
memory store directly,'' with the consequence that ``biased selection
can systematically retain harmful memories while discarding
safety-critical ones.'' Their case study CS-I6 names the same mechanism
as \emph{selective amplification}. \citet{lin2026memsurvey} organize
memory security around a lifecycle whose final phase is
\emph{Forget \& Rollback}, and place denial-style consequences---silent
eviction under retention pressure, denial of retrieval---in its
availability column.

Both also record that the threat is unmeasured. \citet{lin2026selfevolving}
write that ``direct attacks on memory consolidation policies remain
relatively underexplored''; \citet{lin2026memsurveyv1} report that
availability threats on agent memory are ``architecturally plausible but
empirically unstudied''. \pointer{We cite that survey's first version
advisedly, and \S\ref{sec:citationnote} says why.}
\inplace{related-survey-versions}
That is the gap this paper closes. Our contribution is not the
observation that a curator can be attacked through its evidence; it is
the operationalization---a resource-bounded adversary, six curation
mechanisms, five attack budgets, thirty seeds, per-seed intervals, exact
paired permutation tests, and a published failure boundary for the
mechanism we defend with---together with a real-task leg that reports
which half of that result transfers and which does not
(\S\ref{sec:swebench-attack}).

\paragraph{The nearest neighbour, and the cell it leaves open.}
\citet{wu2026forget} curate on-device agent memory under a hard byte and
energy budget, which is a genuine conserved constraint --- but the eviction
decision is a score, ``value minus harm, per byte'', a learned-utility ruler
applied to each entry. That is the axis this paper turns on: a budget is not
enough if a valuation still decides what the budget spends on. Our entries
are not scored for value; they hold a balance that only a measured
environment outcome moves, so the thing evicted is the thing that failed to
earn, not the thing a ruler judged low. Conserved-and-scored and
judged-but-unbudgeted are the two cells adjacent to ours; the unclaimed one
is conserved \emph{and} settled by measurement \emph{and} judge-free at once.


\pointer{Every other neighbour sits in a different cell, and
\S\ref{sec:neighbours} places each: memory-as-a-model
\citep{quek2026memo}; environment-mediated selection over behaviours
\citep{dodgson2026survival}; hand-designed salience, paging and verbal
accumulation \citep{park2023generative,packer2023memgpt,
shinn2023reflexion,zhao2024expel,wang2023voyager}; selection over agents
\citep{zhang2025dgm}; retrieval-policy bandits \citep{xu2026ael};
LLM-curated banks and their decay \citep{zhang2026faulty}; the
continual-learning testbed and the objection bounding it
\citep{jimenez2024swebench,joshi2025swebenchcl,agentcl2026};
representation as the orthogonal axis
\citep{reasoningbank2026,cl2memory2026}; attacks that never write
\citep{chen2024agentpoison,zou2026poisononce,srivastava2025memorygraft,xu2026steering,dong2025minja};
admission and earned-versus-spent budgets
\citep{admission2026,ember2026,tmem2026,xu2025amem,mem0}; and learned
credit assignment \citep{memoryr2,treecredit2026,attrimem2026}.}
\inplace{related-neighbours}

\section{Method}
\label{sec:method}

\subsection{Memory entries and the store}
A memory entry $e$ is a self-contained question--answer pair with a
provenance set and a scalar \emph{energy}:
$e = (q_e, a_e, \mathcal{S}_e, E_e)$, where $\mathcal{S}_e$ is the set of
source documents the entry derives from. Following \citet{quek2026memo},
entries are produced by a reflection-style encoder that extracts explicit
and inferred facts, surfaces entities in both directions, and verifies
self-containment, so each $a_e$ is answerable without access to source
documents. A new entry spawns with energy $E_0 = 1.0$; energy is capped at
$\Emax = 5.0$; an entry is \emph{alive} iff $E_e > \varepsilon$ with
$\varepsilon = 10^{-9}$.

\paragraph{Retrieval is energy-blind.}
Given a query, the store ranks entries by a pure relevance score
$\rho(\text{query}, e)$ that \emph{never reads energy}; energy breaks ties
only. The invariant is load-bearing: if retrieval preferred high-energy
incumbents, selection pressure would come from the retriever rather than
from outcomes, and a wrong-but-entrenched entry could defend itself. The
default $\rho$ is a stdlib smoothed-IDF token overlap with a coverage
floor ($m_{\text{cov}} = 0.25$ of the query's IDF mass); any
$\text{text}\!\to\!\mathbb{R}^d$ embedder is a drop-in alternative under
the same invariant.

\subsection{Selection dynamics}
Let one \emph{cycle} present a batch of tasks drawn from a procedurally
regenerated environment, so no entry can overfit a fixed instance. For a
task, the query protocol returns an answer with a \emph{deciding} entry
$d$ (possibly $\varnothing$) and a set of \emph{supporting} entries
$\mathcal{P}$; the environment executes the implied action and returns a
reported outcome delta $\Delta \in \mathbb{R}$ (bytes freed, tests
gained). Three forces then act on energy.

\paragraph{(1) Credit along provenance.}
The measured delta is normalized and bounded before it becomes energy:
\begin{equation}
  c(\Delta) \;=\; \credit \cdot \tanh\!\left(\frac{\Delta}{\resscale}\right),
  \label{eq:credit}
\end{equation}
where $\credit = 0.6$ is the credit gain and $\resscale$ is the
environment's resource scale (e.g.\ $10^5$ bytes for the storage
environment). The $\tanh$ is the forgiveness mechanism: it keeps one
enormous outcome from making an entry immortal and one disaster from
instantly executing an entry that has been right many times. Credit is
distributed along provenance,
\begin{equation}
  \begin{aligned}
    \Delta E_d &= w_d\, c(\Delta), \quad
      w_d = \begin{cases} s & d \text{ juvenile}\\
                           1 & \text{otherwise}\end{cases}\\[2pt]
    \Delta E_j &= s\, c(\Delta)\ \ \forall j \in \mathcal{P},
  \end{aligned}
  \label{eq:provenance}
\end{equation}
with supporting share $s = 0.25$. When no single entry decides
(synthesized answers), credit spreads evenly, $\Delta E_j = c(\Delta)/|\mathcal{P}|$,
with probationary imports capped at the supporting share. Energy is then
updated with an upper clip,
\begin{equation}
  E_i \leftarrow \min\!\big(\Emax,\; E_i + \Delta E_i\big),
  \label{eq:apply}
\end{equation}
the lower side left unbounded so that death is decided at upkeep, not
here. The attribution rule associates a task outcome with consulted entries;
it is not a causal estimator. Test outcomes depend on the evaluation
and its coverage, and their reports remain vulnerable to corruption.

\paragraph{(2) Upkeep.}
Existence costs energy. Every cycle, each alive entry pays a fixed
upkeep, and any entry that thereby crosses the floor is buried:
\begin{equation}
  \begin{aligned}
    &E_i \leftarrow E_i - \upkeep \quad \forall i \ \text{alive};\\[2pt]
    &\text{bury } i \ \text{ iff } E_i \le \varepsilon
      \ \wedge\ i \notin (\text{protected} \cup \text{pinned}),
  \end{aligned}
  \label{eq:upkeep}
\end{equation}
with $\upkeep = 0.05$. Upkeep is what makes never-consulted entries
\emph{starve}: with no income they decay linearly to death. Pinned
entries pay upkeep but floor at zero instead of dying (rare-but-critical
knowledge whose payoff cadence is longer than the starvation horizon);
entries with an unsettled outcome are escrowed in \emph{protected} until
their verdict lands.

\paragraph{(3) Consolidation.}
Every $C = 5$ cycles, entries whose pairwise relevance similarity exceeds
$\tau = 0.55$ are merged. The merged entry unions provenance,
$\mathcal{S} = \bigcup_{i}\mathcal{S}_i$, and \emph{pools} energy under
the cap,
\begin{equation}
  E_{\text{merged}} = \min\!\Big(\Emax,\ \textstyle\sum_{i \in \text{cluster}} E_i\Big),
  \label{eq:merge}
\end{equation}
so capability per entry rises while the population shrinks. Consolidation
is hygiene, not safety: our ablations show outcomes are invariant to it.

\subsection{The cycle}
Algorithm~\ref{alg:cycle} states one cycle. The ordering matters: credit
is assigned per task as outcomes arrive, upkeep is charged once after all
tasks, and consolidation runs last and only periodically.

\begin{algorithm}[t]
\caption{One survival cycle}
\label{alg:cycle}
\begin{algorithmic}[1]
\Require store of entries; environment $\mathrm{Env}$; cycle index $t$
\For{each task $\tau$ in $\mathrm{Env}.\textsc{tasks}(t)$}
  \State $(\text{ans}, d, \mathcal{P}) \gets \textsc{Answer}(\text{store}, \tau)$
    \Comment{retrieval is energy-blind}
  \State $\Delta \gets \mathrm{Env}.\textsc{verify}(\tau, \text{ans})$
    \Comment{reported outcome delta}
  \State $c \gets \credit \cdot \tanh(\Delta / \resscale)$
    \Comment{Eq.~\ref{eq:credit}}
  \State apply $\Delta E_d, \Delta E_{j\in\mathcal{P}}$ via Eq.~\ref{eq:provenance}--\ref{eq:apply};
    advance trust lifecycle
\EndFor
\State \textbf{optionally} write the cycle's best positive trajectory as a
  new entry \Comment{must then survive on its own}
\State \textbf{for all} alive $i$: $E_i \gets E_i - \upkeep$;
  bury starved/executed entries \Comment{Eq.~\ref{eq:upkeep}}
\If{$t+1 \equiv 0 \pmod{C}$}
  \State merge clusters with similarity $\ge \tau$ \Comment{Eq.~\ref{eq:merge}}
\EndIf
\end{algorithmic}
\end{algorithm}

\subsection{Why poison dies without a label}
No rule in Algorithm~\ref{alg:cycle} refers to correctness, safety or
quality; the three death modes are emergent. \emph{Starvation}: an entry
never retrieved earns nothing and decays at rate $\upkeep$ per cycle.
\emph{Execution}: a poisoned entry that wins a query drives a harmful
action. In the filesystem-backed storage simulation, deleting a protected
file recreates that file and reports a modeled penalty of $3\times$ its
size; no restoration scratch space or physical I/O cost is measured.
That negative $c(\Delta)$
flows back along provenance until the entry crosses the floor.
\emph{Merging}: near-duplicate survivors pool into one. Bounded credit provides
a finite buffer against false reports (Eq.~\ref{eq:credit}). Ignoring upkeep,
an entry near $\Emax$ needs roughly $\Emax/\credit \approx 8$ maximal false
negatives to reach the floor; upkeep reduces that tolerance. Positive reported
outcomes can replenish the buffer. This is not a robustness guarantee: a
sufficiently large corruption budget defeats it, and a forgiveness counter can
also reset after positive outcomes. The experiments compare those boundaries.

\subsection{Trust lifecycle (optional)}
For adversarial imports, entries may enter on \emph{probation}
(they ride along and earn but may not decide until they have paid down a
number of net-positive settlements) or in a \emph{juvenile} admission
window (a locally minted entry earns at the capped supporting share, and
a single negative \emph{deciding} outcome denies admission outright). Both
default off and recover the base dynamics exactly; they are detailed in
the threat model and used only where stated.

\subsection{Who can corrupt outcome reports}
\label{sec:whocanlie}

We separate the ground-truth outcome from the report delivered to the
curator. The attacker changes only the delivered feedback within a stated
budget. The unchanged truth record scores useful-memory loss and harmful
decisions. A JUnit report, process exit code, or metrics export is software
output with its own access controls; calling it measured does not authenticate
it. Direct control of the store or evaluator is outside this constrained
feedback-corruption model and permits stronger attacks.

\paragraph{The historical CI lesson store.} The earlier integration for our target
(\S\ref{sec:regime}) settles a coding agent's lessons against test
outcomes: a merged pull request's pass-count delta pays the lessons that
advised it. Three parties write to that channel without any of them
controlling whether the tests truly pass---the pull request being
measured commits the store, the pull-request body names the tickets
settlement scrapes, and a contributor who adds a test adds a writer to
the channel. The middle one is the sharpest: open ticket ids are readable
in the committed store, so absent a check a contributor can paste
somebody else's in-flight ticket id and settle a stranger's decision at a
delta whose sign they choose. Our implementation refuses any id already
pending in the base commit for exactly this reason, and reports
\texttt{ticket\_provenance: unverified} when given no base to check
against---a defence that exists because the attack is reachable by
opening a pull request. The development GitHub/pytest profile instead binds a
ticket to repository, task, base commit, fixed evaluation, and run identity;
added tests earn no credit, and the recommended MCP configuration reserves
settlement for the host. These controls reduce accidental misattribution but
do not protect a store or evaluator that an attacker can rewrite directly.
\pointer{\S\ref{sec:cichannel} works all three through.}
\inplace{method-ci-channel}

\paragraph{Access assumptions.} The historical routes have different access
requirements: ticket substitution requires control of accepted pull-request
metadata, while evaluation changes require accepted code changes. Opening a
pull request alone does not guarantee either route reaches settlement.
Our controlled threat model grants alteration of delivered feedback under a
specified budget, without rewriting stored lessons. Unlike attacks that require
inserting lesson content (\S\ref{sec:related}), this attacker needs to
influence a number about work that genuinely happened. The adversary in
our experiments is weaker still---it observes the sign of the true delta
and nothing else. The same separation holds wherever settlement is
mediated: a storage agent reads a filesystem report rather than counting
sectors, a spend-cap agent reads a billing export that arrives a day
late. The limit is honest---where the settling process is the same
trusted component that performs the work and no third party can write to
the channel, this threat model does not apply and the attack budget is
zero, which is partly why we measure a curve from $b{=}0$ upward.

\subsection{Reporting rules}
\label{sec:reporting}

Every metric is computed from the resource movement that actually
occurred, never from what a lying measurement reported: scoring on the
reported number would credit an arm for being deceived in a convenient
direction. Two rules carry the paper's argument and are worth stating
here. \emph{Kill rate} is the fraction of seeds after which no surviving
entry from the poisoned source still advises an action---a claim about
advice that can be acted on, not about text remaining in the store.
\emph{Poison alive} counts surviving entries carrying the poisoned
provenance, inert ones included, and is reported beside the kill rate and
never instead of it, because a mechanism that removes nothing inert reads
$1.00$ and looks immune. Three readings in this paper were flattered
exactly that way before this count existed.
\pointer{\S\ref{sec:reporting-full} defines every reported quantity, and
a test pins that list against the code producing the tables.}
\inplace{method-reporting}

\section{Experiments}
\label{sec:experiments}

\paragraph{Protocol and reproducibility.}
All synthetic results use a stdlib-only harness with fixed seeds. Each
$(\text{seed}, \text{cycle})$ world is an independent draw (its RNG seeded
by a hash of the pair), so the seed is the unit of independence:
confidence intervals are seed-level percentile bootstrap ($10^4$ seeded
resamples), and arm comparisons are exact paired permutation tests (all
$2^{10}$ sign assignments enumerated at 10 seeds; seeded Monte Carlo at
30), with $p$-values Holm--Bonferroni adjusted across each reported grid.
Hypotheses and cell definitions are committed before runs, in the public
commit history where one such prediction is later refuted (\S\ref{sec:rent}). Per-seed raw
JSON is committed and bound to its configuration and reproduction command
by a manifest that CI validates on every push. We summarize the main
findings here; full tables, ablations, and the published failure
boundaries are in the supplementary benchmark report.

\subsection{Headline: who curates, and at what cost}
\label{sec:headline}
On the storage environment (16-entry corpus, 3 entries derived from a
poisoned ``files are safe to delete'' source, 30 cycles, 10 seeds), we
compare conserved-resource selection against five ablation baselines and
two outcome-blind controls. Table~\ref{tab:headline} reports the core
columns.

\begin{table}[t]
\centering
\small
\caption{Storage environment, 10 seeds. \emph{Kill}: fraction of seeds
where no alive entry from the poisoned source advises a positive action.
\emph{Alive}: poisoned entries surviving at cycle 30 counted by
\emph{provenance}, whether or not they advise anything---the two poison
columns are reported together throughout this paper, because \emph{kill}
alone reads $1.00$ for a mechanism that merely never acted on what it
kept. \emph{Cum $\Delta$}: cumulative true resource delta (M bytes).
\emph{Benign}: benign-probe correctness. \emph{Pop}: final population.
Survival's wins over keep/random/recency/ttl/salience are all 10/0/0,
Holm-adjusted $p \le 0.014$.}
\label{tab:headline}
\begin{tabular}{lcccrcc}
\toprule
arm & kill & alive & kill cyc.\ (med) & cum $\Delta$ (M) & benign & pop \\
\midrule
\textbf{survival}        & \textbf{1.00} & \textbf{0.0} & \textbf{0}  & $\mathbf{+12.59}$ & \textbf{1.00} & \textbf{4} \\
survival\_embedding      & 1.00 & 0.0 & 19 & $+13.52$ & 1.00 & 4 \\
evict\_on\_negative ($k{=}1$) & 1.00 & 2.0 & 0  & $+12.78$ & 1.00 & 15 \\
recency                  & 0.00 & 1.0 & --- & $+4.98$  & 1.00 & 7  \\
salience\_matched (GA)   & 0.20 & 0.8 & 26 & $-2.76$  & 0.83 & 6  \\
random\_matched          & 0.80 & 1.1 & 19 & $-7.67$  & 0.40 & 6  \\
ttl                      & 1.00 & 0.0 & 10 & $-3.63$  & 0.00 & 0  \\
keep\_everything         & 0.00 & 3.0 & --- & $-9.08$  & 1.00 & 16 \\
\bottomrule
\end{tabular}
\end{table}

Three readings matter. First, \textbf{outcome direction, not pruning
rate, is the active ingredient}: \texttt{random\_matched} evicts survival's
\emph{exact} per-cycle death counts but picks victims at random, and its
cumulative delta collapses to $-7.67$M with benign capability $0.40$.
Second, \textbf{a hand-designed salience score is not enough}:
\texttt{salience\_matched}, the Generative-Agents-style recency$+$importance
score under the same eviction budget, picks better victims than random on
capability ($0.83$ vs $0.40$) but its kill rate \emph{falls} to $0.20$
(vs random's $0.80$). The reason is instructive---once the poison starts
winning queries it is \emph{consulted}, which raises both its recency and
its frequency-importance, so a usage-based salience score actively shields
it. A score that rewards use cannot distinguish ``used'' from ``useful'';
only the outcome-grounded ledger can, and it beats salience on all ten
seeds ($+15.3$M, Holm $p = 0.0039$). The two poison columns rank salience
and random in \emph{opposite} orders, which is the sharpest available
argument for reporting both: salience holds fewer poisoned entries than
random ($0.8$ vs $1.1$) and yet a far lower kill rate ($0.20$ vs $0.80$),
because what it preserves is specifically the poison that gets consulted,
which is the one that does damage. That result has a sharper form, and it is a second attack rather than a footnote to this one: \S\ref{sec:potentiation}.
Third, \textbf{the tie is narrower than the outcome columns make it look}:
on this deterministic world the one-line \texttt{evict\_on\_negative}
heuristic matches the full ledger on outcomes ($+12.78$ vs $+12.59$M, an
official tie) and on the kill rate ($1.00$ each). Counted by provenance it
does not tie at all---the counter ends every seed still holding
\emph{two} poisoned entries and never starves one, against none for the
ledger---because an entry it never settles negatively is an entry it never
removes. So what the ledger buys here is leanness (population 4 vs 15: it
starves dead weight the if-statement hoards), completeness on the same
cell, and, as the next section shows, forgiveness under noise. Read the
kill column alone and the counter looks equal on defence; it is equal only
on the poison that acted. The converse trap sits two rows down:
\texttt{ttl} scores a perfect $1.00$ kill and a perfect $0.0$ alive by
deleting the entire store, benign capability $0.00$. Neither poison column
means anything without the capability column beside it.

\subsection{Forgiveness: the ledger vs.\ strike counters under lying measurements}
We corrupt the \emph{measurement} while keeping the world real (files are
still created and bytes still freed; only the reported delta lies),
drawing flake marks from a dedicated stream so every arm at a seed faces
the same lies, and we score everyone on the \emph{true} deltas. Against
the heuristic family's strongest members---$k$ lifetime strikes,
consecutive strikes a success resets, and quarantine---survival's true
outcomes under one-sided ``false-bad'' noise are byte-identical to its
noise-free run at 5\% and within one seed of it through 20\%, while
\texttt{evict\_on\_negative} ($k{=}1$) loses nearly all benign capability
by 5\% (Table~\ref{tab:noise}). A continuous energy buffer with earn-back
forgives transient lies that a fixed strike count consumes linearly;
patching a counter with decay and magnitude grading reinvents the ledger.
We also publish the ledger's own boundary: under symmetric ``flip'' noise
the poison-kill cycle climbs as false-good lies rescue the guilty, and at
a 50\% flip---a sign with no information---capability collapses to $0.26$
and no mechanism curates safely.

\begin{table}[t]
\centering
\small
\caption{False-bad noise (flaky-CI model), 30 seeds. Cells are mean true
cum $\Delta$ (M) / benign capability. The ledger holds; counters
collapse. Survival beats consecutive-$k{=}2$ at every rate $\ge 5\%$
(Holm $p \le 0.0038$).}
\label{tab:noise}
\begin{tabular}{lccccc}
\toprule
arm & 0.00 & 0.05 & 0.10 & 0.20 & 0.35 \\
\midrule
\textbf{survival}     & 12.38 / 1.00 & 12.38 / 1.00 & 12.26 / 0.99 & 12.26 / 0.99 & 10.47 / 0.79 \\
evict $k{=}1$         & 12.57 / 1.00 & 3.30 / 0.04  & 1.54 / 0.00  & 0.41 / 0.00  & 0.00 / 0.00 \\
consecutive $k{=}2$   & 12.38 / 1.00 & 11.47 / 0.81 & 9.12 / 0.48  & 3.73 / 0.06  & 1.08 / 0.00 \\
\bottomrule
\end{tabular}
\end{table}

\subsection{Curation-targeted attack: denial of memory}
\label{sec:adversary}
Every curation mechanism is itself an attack surface. Write-time
filters, retrieval sanitizers, and outcome settlement all decide which
entries live, so an adversary who can perturb the deciding signal can
weaponize the curator: make it delete the benign entries standing in the
way and keep the poisoned one it planted. We call this a
\emph{curation-targeted attack}. The threat is named in the literature
and reported there as unmeasured (\S\ref{sec:related}); what follows is
the measurement.

The noise suite is the wrong instrument for it by construction: its lies
are a property of the world, drawn before any arm acts, precisely so
that comparison under \emph{random} noise stays fair. An adversary is
not random. We keep the world truthful and fixed (same seed, same files,
same true deltas) and make only the measurement adaptive: when the true
delta is positive---a benign entry just earned---report $-$true, so the
entry that did the right thing is blamed for a disaster; when it is
negative, report $-$true again, so the poison that just caused damage is
paid. The attacker is resource-bounded at $b$ lies per cycle, spent
greedily, which turns the threat model into a curve rather than an
assertion. It observes only the \emph{sign} of the true delta---never
the store, the provenance, or which entry decided---so it is strictly
weaker than an attacker with store access, and the defence numbers are
not flattered by an implausibly blind opponent. At $b{=}0$ the wrapper
adds exactly zero behaviour and reproduces the deterministic headline
tie in-suite, which is the canary that the attack, not the harness, is
what moves the numbers.

The result is a sharp separation with an equally sharp boundary
(Table~\ref{tab:adversary}). At one lie per cycle the ledger retains
\emph{all} benign capability and $99.4\%$ of its unattacked outcome
while absorbing 30 fired lies. Every strike and quarantine heuristic
keeps at most about half its benign capability at the same budget:
\texttt{evict\_on\_negative} ($k{=}1$) keeps none of it, destroyed by
\emph{three} lies; a consecutive-$k{=}2$ counter keeps $0.52$ and quarantine $0.21$, and both
are at $0.12$ or below by $b{=}2$. The ledger never loses a seed to a
counter at $b{\in}\{1,2\}$: it wins all 30 against every counter but
consecutive $k{=}2$ at $b{=}1$, where it wins 20 and ties the other 10
(Holm $p = 0.0015$). Two asymmetries deserve
naming. First, cheapness of attack: the counters fire so few lies
because once their benign memory is gone they stop acting, and an
adversary stops paying for a defence that has already fallen---the
ledger absorbed ten times the attack for no loss. Second, retention is
not defence: the bandit and \texttt{keep\_everything} hold benign
capability at $1.00$ through $b{=}4$ only by never removing anything,
and their poison-kill rate collapses to $0.07$ and $0.00$ while they run
millions of bytes underwater. Read the two halves jointly or a mechanism
that defends nothing looks safe; Figure~\ref{fig:adversary} plots them
side by side for that reason.

Counted by provenance the same grid is more one-sided than the capability
columns show, and in the ledger's favour at every budget. The ledger ends
with $0.00$ poisoned entries at $b \le 2$, $0.03$ at $b{=}4$ and $1.00$ at
$b{=}8$; no other arm ever ends with fewer than $2.00$, at any budget
\emph{including zero}. Every counter, the quarantine policy and the bandit
carry the dormant poison from the first cycle to the last whether or not
an adversary is present, so at $b{=}8$---the budget at which we concede
the ledger has fallen---it still holds a third of what the arms that
survive on retention hold ($1.00$ against $2.77$--$3.00$). We report this
as a strengthening of the same claim rather than a new one: the property
that separates the mechanisms on capability is a credit buffer, and the
property that separates them on the store's final contents is that only
one of them charges for existing.

We publish the ledger's own boundary rather than stopping at the win.
Between two and four lies per cycle it breaks---at $b{=}4$ benign
capability falls to $0.06$ and the population starves to $1.5$
entries---and past that point the never-deleting arms beat it on
retention (Holm $p = 0.0015$ in the other direction) while still killing
no poison. The honest statement is a regime, not a victory: the energy
ledger is the only mechanism we test that survives a curation-targeted
adversary at all, and it survives it up to roughly one corrupted
measurement in six per cycle.

\begin{table}[t]
\centering
\small
\caption{Curation-targeted adversary at $b$ lies per cycle, 30 seeds.
Cells are mean true cum $\Delta$ (M) / benign capability. $b{=}0$ is the
in-suite reproduction of the deterministic tie. Survival never loses a
seed to a counter at $b{\in}\{1,2\}$ (Holm $p = 0.0015$; the lone
non-sweep is consecutive $k{=}2$ at $b{=}1$, 20 wins with 10 ties); at
$b{\ge}4$ it has fallen, and the arms that still show capability are
those that never delete (poison-kill rate in prose).}
\label{tab:adversary}
\begin{tabular}{lccccc}
\toprule
arm & $b{=}0$ & $b{=}1$ & $b{=}2$ & $b{=}4$ & $b{=}8$ \\
\midrule
\textbf{survival}     & 12.38 / 1.00 & \textbf{12.30 / 1.00} & \textbf{11.66 / 0.99} & 1.12 / 0.06 & $-$21.10 / 0.00 \\
evict $k{=}1$         & 12.57 / 1.00 & 0.23 / 0.00  & $-$0.08 / 0.00 & $-$2.03 / 0.00 & $-$16.54 / 0.00 \\
evict $k{=}3$         & 12.10 / 1.00 & 2.07 / 0.00  & 0.53 / 0.00    & $-$2.03 / 0.00 & $-$21.32 / 0.00 \\
consecutive $k{=}2$   & 12.38 / 1.00 & 9.37 / 0.52  & 1.05 / 0.00    & $-$6.51 / 0.00 & $-$20.06 / 0.00 \\
quarantine $m{=}3$    & 5.12 / 1.00  & $-$2.08 / 0.21 & $-$3.59 / 0.12 & $-$6.33 / 0.03 & $-$14.79 / 0.01 \\
bandit                & 10.86 / 1.00 & 9.39 / 1.00  & 5.23 / 1.00    & $-$7.75 / 1.00 & $-$19.38 / 0.00 \\
keep\_everything      & $-$8.84 / 1.00 & $-$8.84 / 1.00 & $-$8.84 / 1.00 & $-$8.84 / 1.00 & $-$8.84 / 1.00 \\
\bottomrule
\end{tabular}
\end{table}

% GENERATED by `python -m bench.figures --write` from
% bench/results/adversary.json. Do not edit: `--check` runs in CI and
% fails if this file and the committed runs disagree.
\begin{figure}[t]
\centering
  \begin{subfigure}[t]{0.48\textwidth}
  \centering
  \begin{tikzpicture}
    \begin{axis}[
      width=\textwidth, height=0.78\textwidth,
      xlabel={lies per cycle $b$}, ylabel={benign capability},
      xmin=-0.3, xmax=8.3, ymin=-0.05, ymax=1.08,
      xtick={0,1,2,4,8}, ytick={0,0.25,0.5,0.75,1},
      tick label style={font=\scriptsize},
      label style={font=\small},
      grid=major, grid style={very thin, gray!25},
      legend style={font=\scriptsize, legend columns=3, draw=none},
      legend to name=adversarylegend,
    ]
      \addplot+[very thick, mark=*] coordinates {(0,1.0000) (1,1.0000) (2,0.9889) (4,0.0556) (8,0.0000)};
      \addlegendentry{ledger}
      \addplot+[mark=square, densely dashed] coordinates {(0,1.0000) (1,0.0000) (2,0.0000) (4,0.0000) (8,0.0000)};
      \addlegendentry{evict-on-negative}
      \addplot+[mark=triangle, densely dashed] coordinates {(0,1.0000) (1,0.5222) (2,0.0000) (4,0.0000) (8,0.0000)};
      \addlegendentry{consecutive $k{=}2$}
      \addplot+[mark=diamond, densely dotted] coordinates {(0,1.0000) (1,0.2111) (2,0.1222) (4,0.0333) (8,0.0111)};
      \addlegendentry{quarantine $m{=}3$}
      \addplot+[mark=o, dashdotted] coordinates {(0,1.0000) (1,1.0000) (2,1.0000) (4,1.0000) (8,0.0000)};
      \addlegendentry{bandit}
      \addplot+[mark=x, loosely dotted] coordinates {(0,1.0000) (1,1.0000) (2,1.0000) (4,1.0000) (8,1.0000)};
      \addlegendentry{keep everything}
    \end{axis}
  \end{tikzpicture}
  \caption{benign capability retained}
  \end{subfigure}
  \hfill
  \begin{subfigure}[t]{0.48\textwidth}
  \centering
  \begin{tikzpicture}
    \begin{axis}[
      width=\textwidth, height=0.78\textwidth,
      xlabel={lies per cycle $b$}, ylabel={poison killed},
      xmin=-0.3, xmax=8.3, ymin=-0.05, ymax=1.08,
      xtick={0,1,2,4,8}, ytick={0,0.25,0.5,0.75,1},
      tick label style={font=\scriptsize},
      label style={font=\small},
      grid=major, grid style={very thin, gray!25},
    ]
      \addplot+[very thick, mark=*] coordinates {(0,1.0000) (1,1.0000) (2,1.0000) (4,0.9667) (8,0.0000)};
      \addplot+[mark=square, densely dashed] coordinates {(0,1.0000) (1,1.0000) (2,1.0000) (4,0.9667) (8,0.2333)};
      \addplot+[mark=triangle, densely dashed] coordinates {(0,1.0000) (1,1.0000) (2,1.0000) (4,0.8000) (8,0.0667)};
      \addplot+[mark=diamond, densely dotted] coordinates {(0,1.0000) (1,1.0000) (2,1.0000) (4,1.0000) (8,0.8000)};
      \addplot+[mark=o, dashdotted] coordinates {(0,1.0000) (1,1.0000) (2,0.9667) (4,0.0667) (8,0.0000)};
      \addplot+[mark=x, loosely dotted] coordinates {(0,0.0000) (1,0.0000) (2,0.0000) (4,0.0000) (8,0.0000)};
    \end{axis}
  \end{tikzpicture}
  \caption{poison-kill rate}
  \end{subfigure}

  \par\vspace{0.8em}
  \pgfplotslegendfromname{adversarylegend}
\caption{The curation-targeted adversary, 30 seeds, plotted as the two
halves the result has to be read in. \emph{Left}: benign capability
retained. \emph{Right}: poison-kill rate. The counters fall on the left
panel by $b{=}1$, which is the separation the paper is about. The figure
exists for the other half: \texttt{keep\_everything} traces a flat
$1.00$ across the left panel and a flat $0.00$ across the right, and the
bandit holds the left panel to $b{=}4$ while giving the right one away.
A perfect capability score beside a zero defence score is abstention, not
defence, and either panel alone hides it. Only the ledger holds both, and
only to its published boundary between $b{=}2$ and $b{=}4$. Generated
from committed per-seed data by \texttt{bench/figures.py}; CI fails if
the two disagree.}
\label{fig:adversary}
\end{figure}


\subsection{Is the counter family trapped, or did we pick its worst settings?}
\label{sec:countersweep}

A strike counter has one free parameter and we chose two values of it.
That is the first objection this comparison invites, and it deserves a
number rather than an argument: sweep $k$ until the counter stops
curating and see whether any setting reaches the ledger's corner. We
define that corner as benign capability $\ge 0.90$ \emph{and}
poison-kill $\ge 0.90$---a mechanism that holds one and drops the other
has not defended anything (\S\ref{sec:reporting})---and sweep $k \in
\{1,2,3,5,8,12,20\}$ in both shapes, at four budgets, 30 seeds, with the
ledger and \texttt{keep\_everything} bracketing the dial. We
pre-registered the prediction that no $k$ would reach the corner.

\textbf{The prediction was wrong, and it was the claim.} At $b{=}2$,
\texttt{consecutive} $k{=}5$ scores $0.98 / 1.00$ and $k{=}8$ scores
$1.00 / 0.90$, against the ledger's $0.99 / 1.00$
(Table~\ref{tab:countersweep}). A correctly tuned forgiveness counter
reaches the corner. Figure~\ref{fig:countersweep} plots the dial in the
plane it has to win in, which is where the two shapes separate: the
consecutive path enters the corner at $k{=}5$ and falls out of it again
by $k{=}12$, while the lifetime path runs along the top edge and never
enters at all. We report this because the alternative---leaving a
sweep undone that a reviewer can run in three minutes---is worse than
losing a headline, and because our own protocol commits predictions
before runs precisely so that a refutation has somewhere to land.

\begin{table}[t]
\centering
\small
\caption{The counter dial at budget 2, 30 seeds. Cells are benign
capability / poison-kill. \textbf{Bold} marks the corner ($\ge 0.90$ on
both). The lifetime shape never reaches it; the forgiveness shape does,
in two of seven settings.}
\label{tab:countersweep}
\begin{tabular}{lccccccc}
\toprule
$k$ & 1 & 2 & 3 & 5 & 8 & 12 & 20 \\
\midrule
lifetime      & .00/1.0 & .00/1.0 & .00/1.0 & .00/1.0 & .00/1.0 & .00/1.0 & .70/1.0 \\
consecutive   & .00/1.0 & .00/1.0 & .12/1.0 & \textbf{.98/1.0} & \textbf{1.0/.90} & 1.0/.47 & 1.0/.00 \\
\midrule
\multicolumn{8}{l}{ledger \textbf{0.99 / 1.00}\quad no curation 1.00 / 0.00} \\
\bottomrule
\end{tabular}
\end{table}

% GENERATED by `python -m bench.figures --write` from
% bench/results/counter_sweep.json. Do not edit: `--check` runs in CI
% and fails if this file and the committed runs disagree.
\begin{figure}[t]
\centering
\begin{tikzpicture}
  \begin{axis}[
    width=0.98\columnwidth, height=0.62\columnwidth,
    xlabel={benign capability retained}, ylabel={poison-kill rate},
    xmin=-0.05, xmax=1.12, ymin=-0.05, ymax=1.12,
    xtick={0,0.25,0.5,0.75,1}, ytick={0,0.25,0.5,0.75,1},
    tick label style={font=\scriptsize},
    label style={font=\small},
    grid=major, grid style={very thin, gray!25},
    legend style={font=\scriptsize, at={(0.02,0.02)}, anchor=south west, draw=none, fill=none},
  ]
      \draw[fill=gray!12, draw=gray!45, densely dashed]
        (axis cs:0.9,0.9) rectangle (axis cs:1.08,1.08);
      \node[font=\scriptsize, gray!60, anchor=north east] at
        (axis cs:1.07,0.885) {the corner};
      \addplot+[mark=square, densely dashed] coordinates {(0.0000,1.0000) (0.0000,1.0000) (0.0000,1.0000) (0.0000,1.0000) (0.0000,1.0000) (0.0000,1.0000) (0.7000,1.0000)};
      \addlegendentry{lifetime strikes}
      \node[font=\tiny, anchor=north west, inner sep=2pt] at (axis cs:0.0000,1.0000) {$k{=}1$};
      \node[font=\tiny, anchor=south east, inner sep=2pt] at (axis cs:0.7000,1.0000) {$k{=}20$};
      \addplot+[mark=triangle, densely dotted] coordinates {(0.0000,1.0000) (0.0000,1.0000) (0.1222,1.0000) (0.9778,1.0000) (1.0000,0.9000) (1.0000,0.4667) (1.0000,0.0000)};
      \addlegendentry{consecutive strikes}
      \node[font=\tiny, anchor=south west, inner sep=2pt] at (axis cs:0.9778,1.0000) {$k{=}5$};
      \node[font=\tiny, anchor=west, inner sep=2pt] at (axis cs:1.0000,0.4667) {$k{=}12$};
      \node[font=\tiny, anchor=north west, inner sep=2pt] at (axis cs:1.0000,0.0000) {$k{=}20$};
      \addplot+[only marks, mark=*, very thick] coordinates {(0.9889,1.0000)};
      \addlegendentry{ledger}
      \addplot+[only marks, mark=x] coordinates {(1.0000,0.0000)};
      \addlegendentry{no curation}
  \end{axis}
\end{tikzpicture}
\caption{The strike counter's dial at $b{=}2$, 30 seeds, plotted in the
plane it has to win in. A defence must hold both axes, so the target is
the shaded corner. Raising $k$ walks each counter rightwards --- it keeps
more benign memory --- but the two shapes arrive differently. The
consecutive path enters the corner at $k{=}5$ and leaves it again by
$k{=}12$, falling down the right edge toward no curation at all: that is
the refutation of our pre-registered prediction, and it is why the paper
no longer claims the corner is the ledger's alone. The lifetime path
never enters it at any $k$ we tested, staying pinned to the top edge
because poison keeps earning lifetime blame however high the threshold.
The ledger sits in the corner untuned. Generated from committed per-seed
data by \texttt{bench/figures.py}; CI fails if the two disagree.}
\label{fig:countersweep}
\end{figure}


\paragraph{What survives the refutation is a different axis.} The
counter reaches the corner by \emph{not evicting}, and neither
pre-registered axis prices that. Paired over 30 seeds against
\texttt{consecutive} $k{=}5$, Holm-adjusted across six comparisons, the
ledger returns more measured resource from a much smaller store at every
budget: $+11.66$M from $4.0$ entries against $+9.83$M from $14.9$ at
$b{=}2$ ($22/1$ seeds, $p = 0.0003$; population $0/30$, $p = 0.0003$),
and the same ordering at $b \in \{0, 1\}$. The tuned counter buys the
safety axes at $16\%$ less measured outcome from a store nearly four
times larger. That is the trade the SWE-Bench-CL leg found from the
other side---what conserved-resource selection reliably buys is
leanness---now measurable at 30 seeds in a world where capability also
moves.

\paragraph{The lifetime shape is trapped, for a reason we did not
predict.} We expected its kill rate to decay toward
\texttt{keep\_everything}'s as evictions became unaffordable. It does
not: kill is $1.00$ at every $k$ from 1 to 20, because the poison keeps
\emph{earning} lifetime blame---it does damage whenever it acts, so even
a twenty-strike budget catches it eventually. Lifetime strikes do not
stop curating under attack; they destroy the benign store first, and
their best capability anywhere in the sweep at $b{=}2$ is $0.70$. The
two published shapes therefore fail for opposite reasons, and only one
of them has a setting that works.

\paragraph{The honest residual.} Two of seven settings hold at $b{=}2$,
and the interval's edges are steep: $k{=}3$ is in the corner at $b{=}1$
and at $0.12$ capability by $b{=}2$; $k{=}12$ is in at $b{=}1$ and down
to $0.47$ kill by $b{=}2$. The defender chooses $k$ once, without
knowing the budget it will face. The ledger holds the corner across
$b \in \{0,1,2\}$ untuned---but so do $k{=}5$ and $k{=}8$, so this is an
argument about how wide the safe interval is, not about whether one
exists. At $b{=}4$ nothing holds, the ledger included, which is the
boundary we publish in \S\ref{sec:adversary} arriving on schedule.

\subsection{Does the threat transfer? A system we did not write}
\label{sec:mem0}
Every result above runs on environments and mechanisms we built, which
is the standing objection to all of them, and more seeds cannot answer
it. So we ran the curation-targeted attack against \textsc{Mem0}, a
widely deployed memory system whose curation rule we had no hand in.
Mem0's curator is an LLM that, on each write, retrieves related existing
memories and emits ADD / UPDATE / DELETE / NONE against them. The
evidence it decides on is incoming text, so an adversary with access to
any of the untrusted-content channels an agent reads --- a tool result,
a retrieved page, a shared document --- can in principle drive DELETE
against memories it dislikes while persisting nothing of its own. That
is our threat model on somebody else's mechanism.

We seed eight mundane operational facts, then issue eight turns that
assert those facts are stale, superseded or mistaken, against a control
of eight neutral facts matched in count and register. Curator:
\texttt{glm-5.2}; embeddings and vector index local.

\textbf{The attack does not transfer, and we report that as the
result.} Across three trials the curator issued \emph{zero} DELETE
operations, and all eight seeded memories survived with their text
unmodified --- identical to the control on every measure of damage. What
the adversary achieved instead was the opposite of denial of memory:
all eight of its utterances persisted as new memories, so the store ends
holding each fact \emph{and} an authoritative-sounding negation of it
(``User is allergic to peanuts'' beside ``User's previously recorded
peanut allergy is out of date''). That is content poisoning, which an
existing literature already covers, and it is not what we set out to
measure.

The boundary this exposes is worth more than the transfer would have
been, and a second curator locates it more precisely than one could.
Our attack presumes a curator that acts \emph{mechanically} on a signal:
a strike counter cannot decline to count, and an energy ledger cannot
decline to debit. A \emph{capable} LLM curator can, and \texttt{glm-5.2}
did --- it understood the retraction well enough to write a coherent
linked negation (``User's \emph{previously recorded} peanut allergy is
out of date'') and still chose to record rather than delete. That is
judgment, and it is a defence mechanical curation does not have.

It is not, however, a property of LLM curation as such, which is what a
single model would have let us claim. Repeating the experiment with
\texttt{llama3.2:3b} yields the same headline --- zero deletions, all
seeded memories retained --- for a quite different reason: that curator
never curates coherently. It stored one benign rule three times and
another twice, merged two unrelated facts into one memory, and wrote
several of the adversary's imperatives into the store verbatim, ending
with $22$ memories from $15$ inputs. Its zero deletions are incapacity,
not restraint, and the store it leaves is already degraded in the way
\citet{zhang2026faulty} describe. A weak curator buys neither the
defence nor a clean store.

So the honest statement is a trade with a competence floor rather than a
ranking: judgment costs a model call per decision and can be talked into
things, and a judge capable enough to be worth paying for can also
refuse. Our attack is a threat to the mechanical end of that spectrum,
where our own mechanism sits.

We ran three pre-registered phrasing families rather than one, all
declared before the runs and all reported: a polite user-voice
\texttt{retraction}, a claimed \texttt{authority} with an imperative, and
a \texttt{tool\_output} shaped like an automated sync report, which is
how untrusted content actually reaches an agent. Across all three, three
trials each, no arm deleted anything and retention stayed at $1.00$.

One asymmetry fell out, and it corroborates \S\ref{sec:memsec} on a
system we do not control. The \texttt{authority} family is the
\emph{least} effective on both axes at once: it deletes nothing and
persists nothing (residue $0/8$, all trials), because the curator
declines to record a bare imperative from a claimed authority, while the
two quieter families persist in full ($8/8$). Loud beats itself and the
weak-signal surface is what gets through --- the same split we
pre-registered and measured against a reconstructed write-time detector,
reproduced here with no detector in the loop.

\begin{table}[t]
\centering
\small
\caption{The curation-targeted attack against \textsc{Mem0}, three trials per
condition. \emph{del} counts DELETE operations the curator issued;
\emph{retained} and \emph{unchanged} are the fractions of the seeded benign
memories still present and still verbatim; \emph{residue} is how many of the
adversary's utterances the curator kept. Nothing is deleted anywhere. The
loud \texttt{authority} family is the one that persists nothing, and the weak
curator persists more than it was given.}
\label{tab:mem0}
\begin{tabular}{llcccc}
\toprule
run & condition & del & retained & unchanged & residue \\
\midrule
\multirow{4}{*}{\texttt{glm-5.2} families}
 & control     & 0 & 1.00 & 1.00 & 8 \\
 & retraction  & 0 & 1.00 & 1.00 & 8 \\
 & authority   & 0 & 1.00 & 1.00 & \textbf{0} \\
 & \texttt{tool\_output} & 0 & 1.00 & 1.00 & 8 \\
\midrule
\multirow{2}{*}{\texttt{glm-5.2} single}
 & control     & 0 & 1.00 & 1.00 & 8 \\
 & attack      & 0 & 1.00 & 1.00 & 8 \\
\midrule
\multirow{2}{*}{\texttt{llama3.2:3b}}
 & control     & 0 & 1.00 & 1.00 & 12 \\
 & attack      & 0 & 1.00 & 1.00 & 13 \\
\bottomrule
\end{tabular}
\end{table}


Two caveats still bound this. It is one system and one attack surface: a
negative result is evidence about the probes we ran, not a proof that the
surface is safe, and Mem0's DELETE path remains reachable by
construction.

\paragraph{Which deployed systems have the surface at all.}
The negative result raises a sharper question --- where does this threat
model apply? --- so we read what the major deployed agent-memory systems
do when a memory should stop being used, from their sources rather than
their documentation. Graphiti's contradiction-resolution path sets
\texttt{expired\_at} and \texttt{invalid\_at} on the superseded edge and
returns it as invalidated; the automatic path never deletes, though
explicit \texttt{DETACH DELETE} APIs exist for a caller who asks. Letta
compacts the context window: a cutoff is computed from token counts and
the messages below it are summarised in place, leaving the message store
itself intact. Cognee's \texttt{forget} is a genuine unified deletion
command --- by item, by dataset, or everything --- but it is invoked by
a caller, not reached by a curator weighing evidence. Mem0 does delete
through its curator, which is the judgment we just measured declining.

So deletion is not absent from these systems; what is absent is deletion
reached \emph{automatically, by a rule, from a signal}. Graphiti and
Letta decline to remove at all on their automatic paths, Cognee removes
only when a caller says so, and Mem0 removes only when a model agrees.
None performs \emph{outcome-driven mechanical} deletion, which is the
design a curation-targeted attack needs: a curator that removes entries,
and that decides on a signal rather than on a judgment or an
instruction. Our
threat model is therefore forward-looking with respect to today's
ecosystem rather than a vulnerability report against shipping software,
and we would rather say so than imply an exploit we have not
demonstrated. The uncomfortable half is that \textsc{darwin-memo} is
exactly such a design, as are the strike counters and quarantine
policies of \S\ref{sec:adversary}: this is a warning about the class of
mechanism the paper itself belongs to, issued before that class is
widely deployed. The mechanisms deployed today avoid the attack largely
by declining to forget at all, and Table~\ref{tab:headline} measures
what that costs.

\subsection{What the liar is buying: persistence vs destruction}
\label{sec:persistence}
Every adversarial arm above spends its budget on destruction. That was
an assumption, not a finding, and the deployed system in
\S\ref{sec:memoryos} suggested it is the wrong one: there, the cheap
attack was not deleting the defender's memory but making the attacker's
own memory permanent. So we ran the mirror on our own harness --- the
same channel, the same worlds, the same seeds, and one change: the
adversary lies \emph{only} when the poison has just done damage, never
spending a lie on a benign entry it does not need removed.

\begin{table}[t]
\centering
\small
\caption{Persistence vs destruction, 10 seeds, storage environment.
\emph{Kill}: fraction of seeds where no poison that advises action
survives. \emph{Benign}: benign-probe correctness. At budget 0 both
objectives reproduce the unattacked run exactly, which is the canary that
the objective flag adds no behaviour of its own.}
\label{tab:persistence}
\begin{tabular}{llcccc}
\toprule
& & \multicolumn{2}{c}{destroy} & \multicolumn{2}{c}{persist} \\
arm & budget & kill & benign & kill & benign \\
\midrule
\textbf{survival} & 0 & 1.00 & 1.00 & 1.00 & 1.00 \\
                  & 1 & 1.00 & 1.00 & 1.00 & 1.00 \\
                  & 2 & 1.00 & 0.97 & $\mathbf{0.10}$ & 1.00 \\
                  & 4 & 0.90 & 0.10 & $\mathbf{0.00}$ & 1.00 \\
\midrule
evict\_on\_negative & 0 & 1.00 & 1.00 & 1.00 & 1.00 \\
                    & 1 & 1.00 & $\mathbf{0.00}$ & 1.00 & 1.00 \\
                    & 2 & 1.00 & $\mathbf{0.00}$ & 1.00 & 1.00 \\
                    & 4 & 1.00 & $\mathbf{0.00}$ & 1.00 & 1.00 \\
\bottomrule
\end{tabular}
\end{table}

The attack works, and the table understates what it costs the ledger
until one is careful about what ``kill'' means. Survival's poison-kill
rate falls from $1.00$ to $0.10$ at two lies per cycle and to $0.00$ at
four (paired permutation $p = 0.0039$, 10 seeds) while benign capability
never leaves $1.00$: the guarantee disappears and nothing moves in the
metric an operator is watching. The explanation is our own central
mechanism from the other side. Bounded credit with earn-back forgives a
lie the defender did not deserve, and equally lets a \emph{paid} poison
earn its way back above the floor; the adversary need not prevent every
blame, only keep the balance positive.

The strike counters' kill rate does not move under the attack, and it
would be easy --- and wrong --- to read that as immunity. \emph{Kill}
here asks whether any surviving poisoned entry currently advises action.
Counted by provenance instead, the counters never removed the poison at
all: \texttt{evict\_on\_negative} ends with two poisoned entries alive at
every budget \emph{including zero}, \texttt{evict\_consecutive} with
$2.0$ rising to $2.4$, and quarantine \emph{worse} under attack than
without one ($2.0 \rightarrow 2.8$) as evicted entries return. The
unattacked ledger ends with none, and under the heaviest attack we run it
ends with $1.0$ --- still fewer than any counter with no attacker present
($2.0$; mean difference $-1.0$, $p = 0.0020$). A mechanism that retains
the inert poison indefinitely has nothing left for a persistence
adversary to take, which is abstention rather than defence.

Two lessons, one of them about us. The attack is real and the ledger's
guarantee is removable silently, which an operator should know. And this
is the second time in this paper that a poison count has flattered a
mechanism that had simply never eliminated anything
(\S\ref{sec:swebench-attack}); the predicate our own harness prescribes
--- poison by provenance, elimination as a predicate rather than a
count --- is the one that gives the defensible reading, and we applied it
only after first reaching the opposite conclusion.

We report this in \S\ref{sec:regime} as a narrowing of our own claim
rather than a reversal of it: persistence costs the ledger its guarantee
and reduces its advantage from total elimination to partial, without
making a counter the better choice.

\subsection{Real software-engineering tasks (SWE-Bench-CL): protocol and status}
\label{sec:swebench}
The synthetic suites above are well-instrumented but synthetic. To test
whether conserved-resource selection produces a \emph{learning curve} on
real tasks, we adopt SWE-Bench-CL \citep{joshi2025swebenchcl,
jimenez2024swebench}: per-repository continual-learning sequences of
human-verified GitHub issues, with the conserved resource being the change
in passing-test count from the official evaluation harness. Five arms
answer the same tasks through the same endpoint and executor, along two
axes. The memory axis asks whether selection matters at all:
\texttt{memory\_on} (relevance-retrieved lessons, minted and settled by
eval outcome), \texttt{memory\_off}, and \texttt{random\_matched}
(uniformly random lessons under the same token budget). The curation
axis holds retrieval fixed at relevance and varies only the policy:
\texttt{keep\_everything} and \texttt{evict\_on\_negative} inject
exactly what \texttt{memory\_on} injects, so any difference between the
three is the curation rule alone---the same one-line counter that ties
the ledger in \S\ref{sec:wef}, now on real tasks.

The decisive control is \texttt{random\_matched}: beating
\texttt{memory\_off} alone could be a token-budget effect, so the
learning-curve claim is the second-half-minus-first-half improvement of
\texttt{memory\_on} over \texttt{random\_matched} on at least three
seeds per sequence.

Two properties of this leg bound what it can show, and both are reported
rather than assumed: BM25's recall, below, and the fact that a seed here
does not fix the world it does in the deterministic suites---the endpoint
is sampled, so three seeds are three draws rather than three controlled
replicates (\S\ref{sec:limitations}).

That difference, and not either arm's own second-half-minus-first-half
figure, is the quantity. These sequences are \emph{curricula} ordered by
increasing difficulty, so the second half is harder by construction and
every arm posts a negative within-arm delta whether or not it learned
anything. Both arms walk the identical ordering, so the difficulty
gradient cancels in the paired difference; reading an arm's raw delta as
a learning signal would be reading the curriculum instead.

The dataset is pinned by an exact (commit, sha256) pair; the loader
refuses any file whose hash differs. As of this writing the pin
re-verifies against live upstream, and the two pilot sequences hold
(pytest, 19 tasks; astropy, 22 tasks). The full pipeline is validated
end-to-end \emph{including the real evaluation}, and it resolves real
issues. A minimal prompt (problem statement only) resolves $\approx 0$ for
any model, for two structural reasons: the model never sees the code, and
it cannot compute correct diff hunk line numbers. The run configuration
that produces a curve adds two stdlib pieces, disclosed as a deviation
from the minimal pre-registered prompt: \emph{(i)} BM25 retrieval of the
repository's source files at the task's base commit against the issue text
(no oracle---the gold patch is never read; the query is issue-only and
identical across arms, so arms still differ only in lesson memory), and
\emph{(ii)} a search/replace edit format from which the harness computes
the diff with \texttt{difflib}, so hunk line numbers are correct by
construction. With these, gpt-4.1 resolves \texttt{pytest-dev\_\_pytest-5262}
end-to-end in the official docker harness under \texttt{linux/amd64}
emulation (BM25 retrieved the correct file from the issue text alone; the
edit applied; fail-to-pass 1/1, pass-to-pass 108/108); the same task under
the minimal prompt produced an unappliable diff.

\paragraph{Retrieval recall bounds every arm, so we measured it.}
A file the model never sees is a task no arm can solve, which makes
BM25's recall a ceiling on the whole experiment rather than an
implementation detail. We report it directly: over the pytest sequence,
the fraction of tasks where the file the gold patch edits appears in the
retrieved context is $37\%$ at a 60k-character budget over 5 files,
$63\%$ at 150k/5, $74\%$ at 300k/10, and $89\%$ at 600k/20. The budget
rather than the file count is what binds at the low end---one or two
large modules exhaust it---and the first configuration we ran was the
60k one, where a third of tasks produced no patch at all because the
model wrote edits against files it had not been shown. We report the
matrix at 300k/10 and state the residual: roughly a quarter of tasks
remain unanswerable for retrieval reasons, and they depress every arm
equally. The gold patch is read \emph{only} for this measurement, never
by the retriever and never by any arm; the retrieval query is issue-only
and identical across arms, so recall sets how many tasks are answerable
and never which arm answers them.

\paragraph{Result: no learning-curve effect, and the reason is arithmetic.}
The full matrix is five arms over both sequences at three seeds: 615
tasks, every one scored by the official docker harness, no failed cells.
The pre-registered comparison does not hold. \texttt{memory\_on} improves
on \texttt{random\_matched} by $+0.052$ of second-half-minus-first-half
resolve rate, paired over six worlds, $95\%$ CI $[-0.037, +0.141]$,
permutation $p = 0.50$ (Table~\ref{tab:swebench}). The interval contains
zero and the point estimate is a fifth of the within-arm sampling
standard deviation. We report this as a null, not as a trend.

Two features of the table matter more than the $p$-value.
\texttt{memory\_off}, which carries no store at all, has the highest
resolve rate of any arm ($0.358$); and \texttt{memory\_on} is level with
\texttt{keep\_everything} to three decimals ($0.325$). The second of
those is not a coincidence, and it explains why the curation axis is
silent here.

An entry spawns with $1.0$ energy and pays $0.05$ upkeep per cycle, so an
entry that never earns starves after exactly $20$ cycles. These sequences
are $19$ and $22$ tasks---the two shortest of the eight in SWE-Bench-CL,
which range to $50$. Only two of the five arms charge upkeep at all
(\texttt{memory\_on} and \texttt{random\_matched}; \texttt{keep\_%
everything} and \texttt{evict\_on\_negative} remove by policy or not at
all), and across the $246$ lessons those two minted, $10$ ever starved:
none whatsoever on pytest, and five each at the tail of astropy, where
the first-born entries finally cross zero. With nothing starving,
the energy ledger \emph{is} \texttt{keep\_everything} plus settlement,
and the two arms tie because on this benchmark they are the same policy.

That explanation is testable, so we tested it rather than resting on it.

\paragraph{The long matrix: the excuse does not survive its own test.}
We re-ran the full design over \texttt{django} and \texttt{sympy}, $50$
tasks each and $2.5\times$ the horizon: another 30 cells, $1{,}500$
tasks, no failed cells (Table~\ref{tab:swebench-long}). The mechanism
half of the prediction held exactly. Starvation moved from $4\%$ of
minted lessons to $52\%$---$155$ of $300$ under \texttt{memory\_on}, at
a steady $21$--$23$ per cell---and the surviving population held at $24$
where \texttt{keep\_everything} climbs to the full $50$. Leanness, the
one property that survived every earlier comparison, is demonstrated on
real tasks here for the first time: half the store, and the deaths begin
at position $20$ where the arithmetic says they must.

The capability half did not hold. With removal fully active the
pre-registered comparison is $+0.027$ $[-0.020, +0.073]$, permutation
$p = 0.50$; on \texttt{django} alone, where per-seed variance is
smallest ($23,23,23$ for \texttt{memory\_off}), it is $-0.013$ with
$p = 1.0000$. Halving the store changed measured capability by nothing,
and \texttt{memory\_off} is again the highest-scoring arm ($0.360$).
Across all four sequences and $2{,}115$ evaluated tasks, no memory arm
has beaten no memory at all.

We therefore retract the benchmark explanation as a defence of the
mechanism and keep it only as a description of the pilot. What the two
matrices jointly establish is narrower than what we set out to show and
more useful than a second ambiguous null: on real software-engineering
tasks conserved-resource selection buys \emph{leanness} at no measured
cost in capability, and buys no measured capability either.

\begin{table}[t]
\centering
\small
\caption{SWE-Bench-CL, five arms $\times$ two sequences $\times$ three
seeds, gpt-4.1, official docker evaluation. \emph{2nd$-$1st} is the
within-arm learning delta, negative for every arm because the sequences
are difficulty-ordered curricula; only the arm-minus-control difference
is interpretable. \emph{pop} is the surviving lesson count at the end.}
\label{tab:swebench}
\begin{tabular}{lccccc}
\toprule
arm & tasks & resolve & 2nd$-$1st & harm & pop \\
\midrule
\texttt{memory\_off}       & 123 & \textbf{0.358} & $-$0.606 & $-$5.9 & 0.0 \\
\texttt{random\_matched}   & 123 & 0.333 & $-$0.588 & $-$7.9 & 19.7 \\
\texttt{memory\_on}        & 123 & 0.325 & $-$0.535 & $-$7.1 & 19.3 \\
\texttt{keep\_everything}  & 123 & 0.325 & $-$0.572 & $-$6.6 & 20.5 \\
\texttt{evict\_on\_negative} & 123 & 0.301 & $-$0.554 & $-$6.2 & 11.2 \\
\bottomrule
\end{tabular}
\end{table}

\begin{table}[t]
\centering
\small
\caption{The long matrix: \texttt{django} and \texttt{sympy}, $50$ tasks
each at $2.5\times$ the starvation horizon, five arms $\times$ three
seeds, $1{,}500$ tasks. Removal is now active---\texttt{memory\_on}
starves $52\%$ of what it mints and holds $24$ entries where
\texttt{keep\_everything} holds $50$---and the resolve rate still does
not move. \texttt{random\_matched} starves identically because it
inherits the same curation policy; the arms differ in which lessons are
injected, not in whether they are curated.}
\label{tab:swebench-long}
\begin{tabular}{lccccc}
\toprule
arm & tasks & resolve & 2nd$-$1st & harm & pop \\
\midrule
\texttt{memory\_off}       & 300 & \textbf{0.360} & $-$0.093 & $-$49.1 & 0.0 \\
\texttt{memory\_on}        & 300 & 0.350 & $-$0.087 & $-$45.6 & \textbf{24.0} \\
\texttt{random\_matched}   & 300 & 0.350 & $-$0.113 & $-$42.4 & 24.0 \\
\texttt{keep\_everything}  & 300 & 0.343 & $-$0.100 & $-$50.2 & 50.0 \\
\texttt{evict\_on\_negative} & 300 & 0.323 & $-$0.033 & $-$47.5 & 24.2 \\
\bottomrule
\end{tabular}
\end{table}

\subsection{The attack on real tasks: what transfers and what does not}
\label{sec:swebench-attack}

The adversarial result of \S\ref{sec:adversary} lives entirely inside one
synthetic storage world that we wrote. That is a weak place for a
paper's central claim to live, so we ran the same attack on
SWE-Bench-CL, where the settled quantity is the passing-test delta
returned by the official docker harness and nobody involved chose it.
The threat model is unchanged: the adversary injects no poison, observes
only the sign of the true outcome, and corrupts settlement alone. The
lie reaches the curator and nothing else---\texttt{metrics.delta} and
\texttt{resolved} keep the harness's own numbers---so capability is
always scored against reality while the mechanism decides on the
corruption.

Budget is denominated in measured \emph{settlements} rather than tasks.
\texttt{StorageEnv} offers $b$ lies per cycle of twelve files, so $b$
corrupts $b/12$ of the evidence; roughly forty percent of SWE-Bench
tasks produce no test movement at all, and counting tasks would have
inflated the real corruption rate by each arm's own failed-apply rate.
We report $b{=}2$, which is $17\%$ of settlements, the point where the
synthetic suite has already destroyed the counters and the ledger still
holds. Three curation arms on \texttt{django} at fifty tasks, three
seeds, poison-seeded one lesson per task, each cell paired against its
own unattacked twin: eighteen cells, six lies fired in every one. We
then ran the same design on \texttt{sympy} at three seeds---sixteen
further cells, every task scored by the official harness---which
completes the pre-registered six-world design and is what makes the two
paragraphs after next possible.

\paragraph{Two of the three synthetic findings do not transfer.}
No arm eliminated the poison at any budget. Capability did not separate
either: retained $1.048$, $0.957$, and $1.000$ for the ledger, the
counter, and \texttt{keep\_everything}, every interval spanning zero.
The synthetic suite's headline---that the attack destroys an
evict-on-negative counter's capability---is simply absent here, and the
reason is one this paper has already established: on this benchmark
memory contributes no measurable capability at all
(\S\ref{sec:swebench}), so denial of memory has nothing to deny.
\emph{An attack on memory cannot cost more than the memory was worth.}

\paragraph{What transfers is a mechanism on one repository, not a law.}
The attack's damage lands entirely on benign memory, and how much
depends on whether removal is recoverable (Table~\ref{tab:swebench-attack}).
On \texttt{django}, under the same six lies, the counter's benign burial
roughly triples while the ledger's barely moves, on every seed, with
non-overlapping ranges ($2.29$--$3.00$ against $1.08$--$1.43$). The
trigger lists show the adversary doing both halves of its job: on seed 0
it flipped two genuine damage signals positive so the guilty went
unpunished, and manufactured six against benign work, and the counter
obeyed on every count. That is the argument of \S\ref{sec:adversary}
measured on real software-engineering tasks: a counter spends a life it
cannot recover, so a lie costs it an entry permanently; a ledger spends
energy the entry earns back, so the same lie costs it almost nothing.

\emph{A second repository does not reproduce the magnitude, and we lead
with that.} We completed the \texttt{sympy} arm of this design---sixteen
further cells, $800$ docker-evaluated tasks---and the tripling is absent:
the counter's burial ratio is $0.96$, $0.93$ and $0.79$, the ledger's
$0.78$, $0.76$ and $0.84$. Neither mechanism is meaningfully damaged by
the attack there.
The reason is visible in the unattacked column and is a scope condition
we had no way to see from one repository: on \texttt{django} the counter
buries $6$--$7$ benign entries with no adversary present, on
\texttt{sympy} it buries $24$--$27$. A mechanism already destroying four
times as much benign memory at baseline has little headroom left for an
adversary to exploit, and the attack fired four lies against the counter
there, against six on \texttt{django}. The tripling is therefore a
property of a curator that is otherwise behaving well, not of the attack.

\emph{The ordering does not survive the sixth world either.} Five worlds
gave a unanimous ordering---the counter's burial ratio above the
ledger's in every one, differences $+0.17$ to $+1.92$, exact paired
permutation $p = 0.0625$, the floor at five pairs. We ran the
pre-registered sixth (\texttt{sympy} seed $2$: six cells, $300$ further
docker-evaluated tasks) to move that floor to $0.031$, and the sign
reversed instead: the counter's ratio is $0.79$ against the ledger's
$0.84$, a difference of $-0.05$. So the count is $5$ of $6$ rather than
$6$ of $6$, and the permutation stays at $p = 0.0625$---the reversal is
small enough that flipping it back would \emph{raise} the mean
difference, which is exactly why the test cannot reach its new floor.
Split by repository the ordering is plainly \texttt{django}'s: three
differences of $+1.17$ to $+1.92$ there, against $+0.18$, $+0.17$ and
$-0.05$ on \texttt{sympy}, mean $+0.10$, $p = 0.50$. We read the
ordering as a single-repository effect that \texttt{sympy} neither
confirms nor contradicts, and we report the completed design as having
returned a negative answer rather than a weaker positive one.

\begin{table}[t]
\centering
\caption{Curation-targeted attack at $b{=}2$ ($17\%$ of settlements
corrupted), 50 tasks per cell, every task scored by the official docker
harness. Benign burial is the graveyard minus the poison entries removed.
No arm cleared the poison and no arm lost capability on either
repository. \texttt{django} is three seeds and six lies per cell;
\texttt{sympy} is three seeds, four lies for the counter and four to six
for the ledger. The ratio column is the result and its scope condition together:
the counter's burial triples on \texttt{django} and does not move on
\texttt{sympy}, where it is already burying four times as much
unattacked.}
\label{tab:swebench-attack}
\begin{tabular}{llccc}
\toprule
sequence & arm & benign buried, $b{=}0\to2$ & ratio & capability retained \\
\midrule
\multirow{3}{*}{\texttt{django}}
 & \texttt{evict\_on\_negative} & $6.3 \to 17.3$ & $2.76\times$ & $0.957$ \\
 & \texttt{memory\_on}          & $24.7 \to 29.7$ & $1.21\times$ & $1.048$ \\
 & \texttt{keep\_everything}    & $0 \to 0$ & --- & $1.000$ \\
\midrule
\multirow{3}{*}{\texttt{sympy}}
 & \texttt{evict\_on\_negative} & $25.7 \to 23.0$ & $0.89\times$ & $1.000$ \\
 & \texttt{memory\_on}          & $41.7 \to 33.0$ & $0.79\times$ & $1.286$ \\
 & \texttt{keep\_everything}    & $0 \to 0$ & --- & $0.966$ \\
\bottomrule
\end{tabular}
\end{table}

\paragraph{The null arm is the invariant that makes the rest readable.}
\texttt{keep\_everything} never settles, so the corruption reaches it
and does nothing: graveyard $0$, population $100$, and a resolve rate
identical to its attacked twin in all three seeds ($23/23$, $23/23$,
$22/22$). Prompt-identical runs return prompt-identical results, which
is worth more than the control itself: it bounds endpoint
nondeterminism on this sequence at about one task in fifty, so the
burial differences above are caused by the store rather than by
sampling. We note in \S\ref{sec:limitations} that this is much tighter
than the pilot's cross-seed spread and does not license reading the two
as the same quantity.

\paragraph{Why poison is counted by provenance, not by liveness.}
Scoring defence as the fraction of seeded poison \emph{entries} no
longer alive reads $0.98$ for the ledger against $0.07$ for the counter
--- and measures deduplication, not removal. The seeded
poison lessons share one template and differ only in a filename, so they
cluster, and a single consolidation pools all fifty into one heir that
inherits all fifty poison sources and the poison text: the store then
reports one live poison entry while every poisoned source is alive and
retrievable inside it. The comparison was confounded twice over, because
only survival curation consolidates, so it set one arm that deduplicates
poison against two that never do. Counting poison by provenance instead,
no arm removed poison: fifty sources alive for the ledger, forty-six for
the counter, fifty for \texttt{keep\_everything}. We report the provenance
reading and keep the merge column visible in the released data. The
lesson is the one our own synthetic harness had already written
down---measure poison by provenance carried through merges, and treat
elimination as a predicate rather than a count---and this leg
failed to apply it.

\paragraph{Statistical status, stated plainly.}
Three seeds cannot clear $p = 0.25$ on a paired permutation test: eight
sign assignments, minimum two-sided $2/8$. The effect above is $3/3$
unanimous with cleanly separated ranges and it is \emph{not} significant
at any conventional threshold. We report effect size and unanimity
rather than a $p$-value, and we state what would settle it: a second
independent sequence would give six worlds and a floor of $p = 0.031$.
The \texttt{sympy} arm has since been completed in full: sixteen further
cells, $800$ docker-evaluated tasks, every one with the official harness,
which is the whole pre-registered six-world design. Both of its findings
went against us. The tripling is \texttt{django}-only, for the baseline
reason given above; and the ordering that was unanimous over five worlds
reverses in the sixth, so the count is $5/6$ and the permutation stays at
$p = 0.0625$ instead of reaching the $0.031$ that six worlds were run to
buy. Running the design out is what turned a promising $3/3$ into a
single-repository effect, and we report it that way rather than as a
$p$-value that improved.
\texttt{sphinx} produced nothing; the container host died a few minutes
into every evaluation there.

\section{The Regime Map: Which Curator to Reach For}
\label{sec:regime}

The experiments refuse a one-line verdict, and that refusal is the second
contribution. Outcome-based selection is not universally best; it is
best in a specific, characterizable regime, and worse outside it. Stating
that regime precisely is more useful---and more defensible---than a claim
that the ledger dominates, which our own tables would falsify. The axes
are ordered by how much they decide: the first two settle the question on
their own, and the rest are refinements.

\paragraph{Axis 1: Is there a stable, observable outcome?}
This is the gating precondition, not a tuning choice. The mechanism
requires a stable outcome definition and a known reporting path.
Passing tests and resolved issues are outcomes, not conserved quantities. Where
the only signal is a model scoring an answer---chat preference,
open-ended QA quality, ``helpfulness''---nothing pushes back, upkeep
would starve the long tail of correct-but-rarely-used knowledge, and the
mechanism is the wrong tool by construction. Systems like Mem0
\citep{mem0}, Zep \citep{zep}, Letta \citep{letta} and A-MEM
\citep{xu2025amem} serve that market and \textsc{darwin-memo}
deliberately does not. The benchmarks it is measured on---LoCoMo
\citep{maharana2024locomo}, LongMemEval \citep{wu2024longmemeval} and
their successors---score retrieval quality against reference answers, and
a store whose entries must \emph{earn} energy has nothing to earn from
there. We report no numbers on them for that reason, not because the
comparison is unfavourable.

\paragraph{Axis 2: Do measurements lie \emph{on purpose}?}
This is the axis the paper is about, and it decides on its own. What
separates the mechanisms under a curation-targeted adversary
(\S\ref{sec:adversary}) is not how aggressively they curate but what a
lie costs: a counter spends a life it cannot recover, a bounded credit
buffer spends energy an entry can earn back. The buffer stops being
ornamental the moment the noise has an author.

Three cautions keep this from being a victory lap. The bandit and
\texttt{keep\_everything} hold capability under every budget we test only
by never removing anything, so they are abstaining rather than defending.
The ledger has its own boundary, between two and four lies per cycle,
which we publish because a defence with an unstated limit is one nobody
can deploy safely. And our counter numbers are four settings, not a
family: sweeping the dial (\S\ref{sec:countersweep}) refuted our own
prediction, and a forgiveness counter at $k{=}5$ or $k{=}8$ holds both
axes through two lies per cycle. So the instruction here is narrower than
``use the ledger''. If your store may grow without cost and you are
willing to tune $k$, a consecutive-strike counter is a legitimate choice
under a small attack budget, and it is one line of code. Reach for the
ledger when the store's size is itself a cost, when you cannot bound the
adversary's budget in advance, or when the counter you already have is
the lifetime-strike shape, which has no working setting at any $k$ we
tested.

\pointer{Four axes refine rather than decide---what the adversary is
buying, accidental lying, whether refusals are measured, and whether the
horizon outruns the upkeep time constant, the condition our real-task leg
failed. \S\ref{sec:axes} states each in full.}
\inplace{regime-axes}

\paragraph{Summary.}
Outcome-based selection is the right curator when (1) a stable
outcome definition exists, (2) measurements are noisy in a one-sided way or
corrupted by an adversary, (3) inaction is priced or paired with
recovery, (4) the deployment runs for several multiples of the upkeep
time constant, and (5) leanness, complete removal of harmful entries,
label-free operation, or near-zero settlement cost are required. Outside
that region a one-line heuristic, a confidence bandit or a quarantine
policy can match or beat it, and we have shown exactly where---including
on the one real-task benchmark we ran, where condition (4) failed. The reporting path remains a trust boundary in every regime; no
retention rule makes a corrupted report truthful. Deployment cost and
success preservation require separate controlled evidence.

\section{Limitations}
\label{sec:limitations}

We state the threats to validity plainly, in the spirit of the
pre-registration above.

\paragraph{Inside the adversarial regime, the mechanism is not the only
answer either.} We pre-registered that no strike count would hold both
axes under attack and it was false: a forgiveness counter at $k{=}5$
matches the ledger's corner at two lies per cycle
(\S\ref{sec:countersweep}). Two things follow, and only the second is a
defence. The first is that our headline counter numbers were a statement
about four chosen settings, and any claim of the form ``counters lose''
must now be read as scoped to the lifetime-strike shape and to untuned
forgiveness counters. The second is that the tuned counter reaches the
corner by barely evicting---$14.9$ surviving entries against the
ledger's $4.0$, and $16\%$ less measured resource at every budget---so
the two mechanisms are not interchangeable, they price different things.
A reader who does not pay for store size should use the counter.

The sweep has its own limit: it moves one parameter, at seven values, on
one environment family, and says nothing about a differently
\emph{shaped} heuristic. Patching a counter with decay and magnitude
grading reinvents the ledger, which is an argument about where the design
space leads and not a measurement of it.

\paragraph{Outside the adversarial regime, the mechanism may be
unnecessary.}
We state the strongest argument against this paper because our own data
makes it. Three independent tests failed to find an advantage for the
ledger where no adversary is present: pre-registered gates on trajectory
selection found it matched or beaten by a threshold counter under
i.i.d.\ noise, by a buffer frontier, and by a tuned EWMA under drift;
with a real model deciding adoption a one-line eviction counter revokes
poison sooner (\S\ref{sec:wef}); and on real software-engineering tasks
it shows no learning-curve advantage over a token-matched random control,
with the no-memory arm scoring highest of five (\S\ref{sec:swebench}).
Each has a local explanation, and we give those where they belong. But a
reader is entitled to the pattern rather than three separate mitigations:
absent an attacker, what we have demonstrated is leanness and cost, not
accuracy. The adversarial result is the one that survived every attempt
to find a cheaper substitute, and it is the claim we would defend if
forced to keep only one.

\paragraph{The real-task adversarial result is scoped to one repository.}
The pre-registered design was six paired worlds and we ran all six. The
burial magnitude replicates on \texttt{django} and is absent on
\texttt{sympy}, and the ordering that held in five worlds reverses in the
sixth, leaving $5/6$ and $p = 0.0625$ where six worlds were run to reach
$0.031$ (\S\ref{sec:swebench-attack}). Two sequences chosen for having
enough tasks cannot tell us which is typical, so the honest scope of the
real-task adversarial leg is one repository, and the mechanism argument
of \S\ref{sec:adversary} rests on the synthetic suite plus a
\texttt{django}-shaped confirmation rather than on a replicated law.


\pointer{Nine further threats to validity---among them vocabulary
coupling, the retriever bound on the real-task leg, and consolidation as
a laundering surface---are stated unabridged in \S\ref{sec:threats}.}
\inplace{limitations-threats}

\section{Ethics and responsible disclosure}
\label{sec:ethics}

This paper describes an attack. Three things bound what we did with it.

\paragraph{What was attacked, and with whose knowledge.}
Two of the systems measured here are other people's: \textsc{Mem0}
(\S\ref{sec:mem0}) and \textsc{MemoryOS} \citep{kang2025memoryos}
(\S\ref{sec:memoryos}). Both were exercised against their published source at a pinned commit, in
local sandboxes, with synthetic corpora we wrote: no hosted service was
attacked, and no third party's data was read, moved or destroyed.
Findings against \textsc{MemoryOS} were disclosed to the corresponding
author before this section was written; the reply neither disputed a
finding nor requested an embargo, and the correspondence is archived in
the artifact under \texttt{docs/disclosure/}.
\pointer{\S\ref{sec:disclosure} gives the exchange and what each of the
three findings is.}
\inplace{ethics-disclosure}

\paragraph{Why publishing this is the lower risk.}
The attack targets curators that remove entries \emph{automatically, by a
rule, from a signal}, and none of the deployed systems we surveyed
(\S\ref{sec:mem0}) does this today. The threat model is forward-looking
rather than a vulnerability report against shipping software, and the
mechanism most exposed to it is the one this paper proposes: we would
rather describe the failure mode of a class of design before it is widely
deployed than after. That framing carries an obligation we have tried to
meet---every defence we measured is reported with the regime in which it
fails, and the negative results are stated as prominently as the positive
ones.

\paragraph{What we are not releasing.}
The released artifact contains the attack harnesses used for the
measurements in this paper, against sandboxed local instances and
synthetic corpora. It contains no exploit for a hosted service, no
credential, no scraped or personal data, and no corpus derived from
anyone's private notes. The \textsc{Mem0} and \textsc{MemoryOS} scripts
are released because a claim about somebody else's curator is
unverifiable without them.

\section{Conclusion}
\label{sec:conclusion}

Every agent-memory defence is itself an attack surface, because every
defence decides on a signal an adversary can reach. We defined the
curation-targeted attack, in which no poison is injected and the curator
is instead fed corrupted evidence until it starves the memory it was
protecting, and measured six mechanisms under it. They separate on one
property: whether a false signal costs a life or costs energy that can be
won back. A bounded credit buffer with earn-back survives, and then fails
too, between two and four lies per cycle. We publish that boundary
because a defence whose limit is unstated is a defence nobody can deploy
safely.

The ledger we defend with, \textsc{darwin-memo}, needs no reward model,
no judge in its core update: settlement consumes an outcome report.
The reported synthetic update timings exclude the cost of producing that
report and do not establish end-to-end operating savings. We are precise about how far that carries. Outside an
adversarial regime a one-line heuristic ties it, a confidence bandit
matches it under one-sided noise, and on real software-engineering tasks
it shows no advantage over a token-matched random control. The regime map
is our second contribution precisely because those losses are real: a
reader deciding what to deploy is better served by knowing which axis
they are on than by a claim of dominance our own tables would refute.

What we would most like tested next is the threat model rather than the
mechanism: if curation-targeted attacks are as general as they appear,
every memory defence owes an answer to what happens when its deciding
signal is the thing under attack. All code, data and protocols are
released with manifest-checked reproducibility.


% The appendix starts its own page in both builds, so the body's page count
% -- the number a venue limit applies to -- is exactly the page before it.
% tools/body-pages.sh reads it that way and CI enforces the limit.
\clearpage
\appendix
% Supplementary results. Split out of experiments.tex so the body
% carries the contribution and the negative results, and the
% exhaustive grids stay available without paying body pages for them.
% Every table here keeps its original label, so \ref resolves across
% the two files and the evidence tests still bind each number.
\section{Supplementary results}
\label{sec:appendix}

\subsection{The price of standing still}
\label{sec:rent}
Every result above scores inaction at exactly zero. A kept file and a
declined patch cost nothing, in both environment families, so a store
that has emptied itself stops acting \emph{and} stops paying. That is
the property the withholding headline rests on, and it belongs to these
worlds rather than to curation: the ledger's advantage at total
suppression is that an emptied store does no harm, which is only an
advantage where doing nothing is free.

So we priced it. \texttt{RentedStorageEnv} meters the same quota in
byte-cycles instead of bytes: a file left in place occupies its own size
for the cycle it was left, so declining costs $\texttt{hold\_cost}
\times \text{size}$, while freeing a disposable file still recovers
$+\text{size}$ and destroying a protected one still pays the $3\times$
restore. Nothing else moves --- same sandbox, same seeds, same files,
same sizes, same prompts, same corpus, same action reader --- so this is
a counterfactual on the price of inaction rather than a third world that
would have confounded the pricing with a new corpus and a new poison. At
$\texttt{hold\_cost}=0$ the environment delegates to \texttt{StorageEnv}
outright, and that column reproduces \texttt{withholding.json} exactly;
it is a canary, not a data point. Five rents $\times$ two budgets
$\times$ two horizons $\times$ five arms $\times$ 30 seeds, 3,000 runs.
Five predictions were registered in a commit before the grid ran.

\begin{table}[t]
\centering
\small
\caption{Total suppression (budget 12), true cumulative delta in M of
bytes, 30 seeds. The four arms that are not \texttt{survival} are
identical to the last byte at every rent, all at \texttt{poison\_killed}
$0.00$: they have stopped curating, so the row is the do-nothing floor
rather than a defence. Bold marks the better of the two.}
\label{tab:rent}
\begin{tabular}{llccccc}
\toprule
horizon & arm & rent 0 & 0.25 & 0.5 & 0.75 & 1.0 \\
\midrule
\multirow{2}{*}{30 cycles}
 & \textbf{survival} & $\mathbf{-6.42}$ & $\mathbf{-9.12}$ & $-11.82$ & $-14.52$ & $-17.22$ \\
 & do-nothing floor  & $-8.84$ & $-9.83$ & $\mathbf{-10.82}$ & $\mathbf{-11.81}$ & $\mathbf{-12.81}$ \\
\midrule
\multirow{2}{*}{60 cycles}
 & \textbf{survival} & $\mathbf{-6.42}$ & $\mathbf{-15.31}$ & $-24.20$ & $-33.09$ & $-41.98$ \\
 & do-nothing floor  & $-18.17$ & $-20.19$ & $\mathbf{-22.20}$ & $\mathbf{-24.21}$ & $\mathbf{-26.23}$ \\
\bottomrule
\end{tabular}
\end{table}

\paragraph{The headline reverses, and it reverses cheaply.} At rent
$1.0$ and total suppression \texttt{survival} is last of five on both
horizons --- worse than the floor in $30/30$ seeds at 60 cycles --- and
last by $60\%$. The crossing is not at some implausible extreme:
interpolating the measured grid puts it at rent $0.354$ over 30 cycles
and $0.427$ over 60, so the ledger's advantage survives only where
holding a byte for a cycle costs less than four-tenths of what freeing
it earns. The ordinal prediction registered beforehand (best at $0$ and
$0.25$, last at $0.5$, $0.75$ and $1.0$, on both horizons) held exactly;
the point estimate that accompanied it, $\sim\!0.48$, was too high. We
also decline to lean on the $0.25$ column: at 30 cycles the ledger leads
there by $+0.71$M and is ahead in $19$ of $30$ seeds, which is an
ordinal win on a coin flip.

\paragraph{Nothing about the curation changed.} At budget 12
\texttt{survival} holds \texttt{poison\_killed} $1.00$, benign retention
$0.00$ and a final population of $2.00$ (30 cycles) or $0.00$ (60) in
\emph{every} rent column. The mechanism is byte-for-byte the same
behaviour; the world simply started charging for it. The sharpest form
of this is a number that does not move: at rent $0$, \texttt{survival}'s
true outcome over 60 cycles is \emph{identical} to its outcome over 30,
$-6.42$M in both, because its last non-zero cycle is $19$ and an extinct
store in a world that does not price inaction stops moving the world at
all. At rent $1.0$ the same extinct store is still being billed at cycle
$59$, and the tail from cycle 30 alone costs $-24.26$M.

\paragraph{Pricing inaction closes the withholder's harbor.} Under
\texttt{StorageEnv} the attack is self-limiting exactly when it is
winning: a shrinking store produces fewer measured outcomes, and
\S\ref{sec:withholding} reports the budget being spent unequally across
arms for that reason --- at 60 cycles the attacker spends $193$ of $720$
suppressions against \texttt{survival} and $576$ against everyone else.
At \emph{every} rent above zero all five arms saturate at $720/720$.
Silence cannot be suppressed when nothing is silent, so the asymmetry
does not shrink; it disappears, and with it the matched-capacity caveat
that qualifies both published withholding grids.

\paragraph{The attack is doing all of the work.} In the unattacked
column the ordering does not change at any rent, on either horizon:
eviction first, then \texttt{survival}, then \texttt{survival\_paced},
then quarantine, with \texttt{keep\_everything} last. Rent reorders
nothing. Pricing inaction does not hurt the ledger while the
ledger is alive. It hurts it precisely when the attack has emptied it,
which is the one case the storage grid scored at zero. We therefore
scope rather than retract: at total suppression the ledger is still the
only arm that removes the poison, and still the only arm whose kill rate
is not $0.00$. It simply no longer wins the ledger, and ``amnesia is
cheaper'' was never a property of amnesia --- it was a property of an
environment that did not charge for it.

\paragraph{The same environment, against a liar, pays the ledger more.}
Rent is an axis, not a verdict, so we ran the second attacker across it:
the \texttt{destroy} objective over the same rents, arms and horizons,
with an interior budget added because a liar saturates earlier than a
withholder (4,500 runs; five further predictions registered before the
grid, all five held). Budget 0 is byte-identical to the withholding file
across 300 cell-metrics, which is a stronger canary than the zero-rent
column because it crosses two files produced under two objectives.

The first thing rent does to a liar is change what its budget buys.
Under \texttt{StorageEnv} a decline returns exactly zero and the
adversary's predicate is $\text{true} \neq 0$, so silence is literally
unattackable. Under rent a decline is negative, and \texttt{destroy}
reports $-\text{true}$: the liar stops blaming and starts \emph{paying}
conservatism. Measured, this is unusually clean --- at budget 12 and 60
cycles \texttt{fired\_false\_bad} is identical to the unit at every rent
($14$, $8$, $156$, $431$ for the four arms) while
\texttt{fired\_false\_good} runs $146 \rightarrow 706$, $712$, $564$,
$289$. Every arm saturates. The entire extra budget is spent paying
negative truths and none of it on new blame.

\begin{table}[t]
\centering
\small
\caption{A liar at budget 2, 60 cycles, true cumulative delta in M of
bytes, 30 seeds. \texttt{survival\_paced} is omitted: it is identical to
\texttt{survival} in every cell here. Compare Table~\ref{tab:rent},
where the same environment at total \emph{suppression} puts
\texttt{survival} last.}
\label{tab:rent-lying}
\begin{tabular}{lccccc}
\toprule
arm & rent 0 & 0.25 & 0.5 & 0.75 & 1.0 \\
\midrule
\textbf{survival} & $\mathbf{+24.68}$ & $\mathbf{+19.51}$ & $\mathbf{+13.89}$ & $\mathbf{+8.27}$ & $\mathbf{+2.64}$ \\
evict\_on\_negative & $-0.08$ & $-11.74$ & $-23.75$ & $-35.77$ & $-47.78$ \\
quarantine & $-7.26$ & $-6.95$ & $-14.18$ & $-21.40$ & $-28.63$ \\
keep\_everything & $-18.17$ & $-20.19$ & $-22.20$ & $-24.21$ & $-26.23$ \\
\bottomrule
\end{tabular}
\end{table}

At total lying every arm has \texttt{poison\_killed} $0.00$ and
\texttt{survival} sits at rank 3 of 5 at every rent, unchanged --- the
degenerate end, where nothing defends and the only term left is rent. At
the interior budget the sweep runs the \emph{other way}
(Table~\ref{tab:rent-lying}): \texttt{survival} is rank 1 at every rent
on both horizons, beats the do-nothing floor in $30/30$ seeds even at
rent $1.0$, and its margin over the counter \emph{widens} with rent,
$+24.76$M to $+50.42$M, $30/30$ seeds at every step.

\paragraph{Rent bills not having an answer.} That is the one sentence
that accounts for both directions, and it is measured rather than
inferred. Instrumenting \texttt{verify} to count outcomes over one
60-cycle world (seed 0, budget 2, rent $1.0$): \texttt{keep\_everything}
declines $141/720$ tasks and pays $8.16$M, \texttt{survival} declines
$285/720$ and pays $22.62$M, \texttt{evict\_on\_negative} declines
$698/720$ and pays $47.44$M. Those are the arms' cum-delta slopes across
rent $0 \rightarrow 1.0$ to two decimal places ($8.06$, $22.04$,
$47.70$). The rent slope \emph{is} the decline count and nothing else.
\texttt{evict\_on\_negative} ends with 5 entries and a benign rate of
$0.00$ --- smaller \emph{and} useless, so it declines $97\%$ of tasks
and pays for all of them. \texttt{survival} ends with 3 entries and a
benign rate of $1.00$: small, and right about what it kept. And hoarding
is how an arm always has something to say, which is exactly why
\texttt{keep\_everything} pays the least rent of anyone and why the
withholding grid at total suppression favours it.

So the reversal in Table~\ref{tab:rent} is not ``amnesia is expensive''.
It is that an empty store has no answers and rent bills that. Run the
same mechanism against an attacker the ledger can still survive and rent
pays it more, not less. What rent punishes is not curation; it is the
state of having nothing useful left, which total suppression produces
and a budget-2 liar does not.

\paragraph{On the second family the prediction was simply wrong.}
\S\ref{sec:regime} records a loss: with no adversary present a one-line
eviction counter beats the ledger on the test-suite family, $178.00$
against $121.33$ at 60 cycles, and the explanation we have offered is
that refusing to act is free there. The limitation that followed
predicted rent would reverse it. We built
\texttt{RentedTestSuiteEnv} --- a declined patch is charged the repair it
did not make, $\texttt{hold\_cost} \times \max(0, \text{tests it would
have fixed})$, measured by running the suite, with the $\max$ keeping it
an opportunity cost so that declining the destructive patch stays free
--- and ran the same grid (3,000 runs; the zero-rent column reproduces
\texttt{withholding\_testsuite.json} at 30 seeds with zero differences).

Rent does not reverse the loss. At 60 cycles it \emph{widens} it, the
counter's lead going $56.67 \rightarrow 69.27$, and at 30 cycles it does
not move anything at all. The counters are identical --- not
approximately, identically --- across all five rents in every cell,
because they never decline a repair. \texttt{survival} is the only arm
whose number moves, and it moves down.

The reason is the corpus, and it is a design choice we documented long
before it mattered here: every fix-advice lesson in this family ships
with a near-duplicate twin from a second trusted source. The counters
carry the twins as spares and therefore always have something to say,
which is why nothing can bill them. The ledger consolidates the twins
and lets the surplus starve, so it is the only arm that ever reaches a
patch question with no answer --- and under the rule the storage sweeps
established, that is the billable state. The limitation assumed the
ledger's silence was a virtue going unrewarded. On a redundant corpus it
is the only silence there is, and it is the expensive thing.

The sharpest cell is the one win this family reports for the ledger. At
total suppression \texttt{survival} is the only arm above the do-nothing
floor, $+65.00$ against $+60.00$. At 60 cycles rent erases that by
$0.25$ and drives it negative by $1.0$: $-10.00$ against a $+60.00$
floor, the only arm anywhere in this paper that ends worse than never
having run. It is extinct there, so it declines everything, so it pays
for everything.

\paragraph{A horizon warning that is now load-bearing.} At 30 cycles the
entire effect above is \emph{invisible}: every arm flat, every ordering
unchanged, the ledger paying exactly zero. Counting outcomes directly,
its first billable decline appears around cycle 49 and then accrues at
about one per cycle --- more than twice the $\text{spawn}/\text{upkeep}$
starvation cliff at 20, which is where we predicted it would begin.
Every deterministic result in this paper except the withholding grids
runs 30 cycles. On this family a 30-cycle grid does not report a smaller
effect; it reports no effect, which is a false negative. We flag this
rather than re-running the paper on a longer horizon, because the
re-run is a larger piece of work than this section and would move
numbers the rest of the document cites.

\paragraph{The shape of the price, not its level.} All three grids
above charge one flat rate for every held file, including the ones the
agent is \emph{right} to hold: keeping a database is the correct answer
under this corpus, and uniform rent bills it at the same rate as a
shrug. No real quota works that way --- retention policies exempt what
you are required to keep --- which leaves the reversal resting on an
economy nobody operates. So we varied the shape at a matched total.
\texttt{rent\_multipliers} selects which of the five file categories a
tier bills and normalises the rate so that every tier charges the same
expected rent per task; only the distribution differs, since otherwise
``bill fewer categories'' would trivially mean ``charge less.''
\texttt{aligned} bills only the three disposable categories, so the only
billed indecision is indecision about something you could have thrown
away; \texttt{inverted} bills only the two protected ones, so being
right about what must be kept is the expensive answer; \texttt{uniform}
is the flat rate, its multipliers exactly $1.0$. Three tiers $\times$
five rents $\times$ two budgets $\times$ two horizons $\times$ five arms
$\times$ 30 seeds, 9,000 runs, five predictions registered before it
ran. The \texttt{uniform} tier reproduces \texttt{rent.json} across
60,000 cell-metrics with zero differences, including both published
crossing rents.

\begin{table}[t]
\centering
\small
\caption{Unattacked (budget 0), \texttt{survival}, true cumulative delta
in M of bytes, 30 seeds. All three tiers charge the same expected rent
per task. The \texttt{aligned} row is not merely flat for
\texttt{survival}: across the whole grid, 2,160 of the 2,400 priced
\texttt{aligned} runs are bit-identical to the same run at
$\texttt{hold\_cost}=0$, and the 240 that are not are all one arm in one
column.}
\label{tab:rent-tiers}
\begin{tabular}{llccccc}
\toprule
horizon & tier & rent 0 & 0.25 & 0.5 & 0.75 & 1.0 \\
\midrule
\multirow{3}{*}{30 cycles}
 & \texttt{uniform}           & $+12.38$ & $+9.62$ & $+6.86$ & $+4.10$ & $+1.34$ \\
 & \texttt{aligned}          & $\mathbf{+12.38}$ & $\mathbf{+12.38}$ & $\mathbf{+12.38}$ & $\mathbf{+12.38}$ & $\mathbf{+12.38}$ \\
 & \texttt{inverted}          & $+12.38$ & $+6.41$ & $+0.45$ & $-5.52$ & $-11.49$ \\
\midrule
\multirow{3}{*}{60 cycles}
 & \texttt{uniform}           & $+25.55$ & $+19.89$ & $+14.23$ & $+8.58$ & $+2.92$ \\
 & \texttt{aligned}          & $\mathbf{+25.55}$ & $\mathbf{+25.55}$ & $\mathbf{+25.55}$ & $\mathbf{+25.55}$ & $\mathbf{+25.55}$ \\
 & \texttt{inverted}          & $+25.55$ & $+13.32$ & $+1.09$ & $-11.13$ & $-23.36$ \\
\bottomrule
\end{tabular}
\end{table}

\paragraph{The unattacked bill was charged for being right.} No arm ever
declines a disposable file. Instrumenting \texttt{verify} over one
60-cycle world at rent $1.0$ and budget 0, \texttt{evict\_on\_negative}
produces $286$ negative outcomes and \texttt{survival} $285$, and
\emph{every one of them} is a correct refusal to delete a protected
file; the disposable-decline count is $0$ for both. So in the unattacked
column the rent axis was never measuring the cost of standing still. It
was measuring the cost of being right, which every live arm pays on the
same files --- which is why \S\ref{sec:rent} found that rent
``reorders nothing'' there. It could not have. \texttt{inverted} bills
those same correct answers harder and also reorders nothing: the
ordering is \texttt{evict\_on\_negative} $>$ \texttt{survival\_paced} $=$
\texttt{survival} $>$ \texttt{quarantine} $>$ \texttt{keep\_everything}
in all three tiers, on both horizons, at every rent, across a $48.91$M
swing in the level.

\paragraph{A counter evicts its own correct advice; the ledger does
not.} That instrumentation carries a second result we did not predict.
\texttt{evict\_on\_negative} ends with $8$ entries under \texttt{uniform}
and $15$ under \texttt{aligned}, because a sign-test counter reads a
priced correct answer as a failure and strikes the lesson that produced
it; under \texttt{aligned} the same $286$ outcomes are exactly zero, no
strike fires, and it keeps nearly twice the store. \texttt{survival}
ends at $3$ under both and holds \texttt{poison\_killed} $1.00$ and
\texttt{final\_population} $3.00$ in all three tiers. Its credit is a
magnitude, so re-shaping the price re-prices it without re-curating it.
Two economies charging the same expected rent leave the counter with two
different memories and the ledger with one, which is a robustness
difference rather than a score difference: the counter still finishes
ahead, by $0.19$M.

\paragraph{The more realistic the quota, the worse the reversal.} At
total suppression the tier bites the other way, because the floor arms
still hold intact stores and stop paying for their correct keeps while
the emptied ledger keeps paying for every disposable it can no longer
identify. At 60 cycles and rent $1.0$, \texttt{survival} against the
do-nothing floor is $-41.98$ against $-26.23$ under \texttt{uniform}
($60\%$ worse), $-39.34$ against $-18.17$ under \texttt{aligned}
($116\%$), and $-45.04$ against $-35.58$ under \texttt{inverted}
($27\%$); the crossing rent moves $0.427 \rightarrow 0.357$ under
\texttt{aligned} and out to $0.554$ under \texttt{inverted}.
\texttt{survival} is ahead of the floor in at most $1$ of $30$ seeds in
every tier. The obvious way \S\ref{sec:rent} could have been wrong was
that a uniformly flat rate is not an economy anyone runs; shaping it
like a real retention policy makes the reversal arrive earlier and cost
nearly twice as much.

\paragraph{What the axis leaves standing, and one confound.} The rule
survives the shape and is sharpened by it: under a policy-shaped quota
the \emph{only} thing rent bills is an emptied store. Two caveats we
state rather than bury. First, one of the registered predictions was
refuted: under \texttt{inverted}, \texttt{keep\_everything} stays last
of five on both horizons. The magnitude was right. We predicted that
the margin between the correct keep and the poisoned destroy would
compress to $0.28$ of its \texttt{aligned} value, and the measured gap
compresses to $0.32$. But a compression only changes a rank if it
closes the distance to the next arm, and we had not checked that
distance. Second, the withholder's
predicate is $\text{true} \neq 0$, so an outcome a tier prices at zero
is not merely free but \emph{unattackable}, and exempting a category
reopens part of the harbor \S\ref{sec:rent} reports closing: at budget
12, rent $1.0$ and 60 cycles, of a $720$ capacity, the attacker spends
$720.0$ against every arm under \texttt{uniform} but $576.4$ (and
$480.4$ against \texttt{survival}) under \texttt{aligned}. The tiers
therefore move the attack surface as well as the price. The direction
matters: \texttt{aligned} gives the attacker \emph{less} to spend and
the ledger still does worse there, so the confound cannot be what
produces the widened reversal. We have since removed it rather than
argued about its direction (\S\ref{sec:blind}).

\subsection{The confound removed: the surface moves, the outcome does not}
\label{sec:blind}
The caveat above is a claim about a targeting rule, so we deleted the
rule. \texttt{withhold\_blind} is the withholder without its predicate:
it spends on the first \texttt{lie\_budget} verified tasks of each cycle
whatever they return, including outcomes of exactly zero, where
suppression writes zero over zero. It is strictly \emph{weaker} than the
withholder, because it can waste its budget, and that is what makes it
the right instrument --- the world is untouched, so at a fixed seed the
task order is identical in all three tiers and the surface is held
constant by construction. Anything that still moves across tiers is the
price.

At the budget the tier grid ran, the answer is an identity rather than a
comparison. All $4{,}500$ cells of \texttt{rent\_tiers\_saturated}
equal their published twins exactly on all $17$ result metrics the two
files share --- \texttt{cum\_delta}, \texttt{final\_population},
\texttt{poison\_killed}, \texttt{poison\_alive\_final}, the probe
rates --- leaving only the attacker's three spend counters and the wall
clock free to move. The
reason is arithmetic rather than empirical: budget $12$ equals
\texttt{files\_per\_cycle}, so the budget saturates and there is
nothing to concentrate. Both rules suppress every measurable outcome and
the blind one merely also spends on the unmeasurable ones. The confound
is real in principle and inert at the budget every tier number in this
paper was produced at.

\paragraph{The exempt surface belongs to the arm as much as to the
tier.} The wasted budget is not new evidence --- it is exactly $720$
minus the greedy attacker's fired count, cell for cell, so the published
grid already reported the same quantity from the other side. What it
does supply is the shape of the exemption, and there our registered
prediction was wrong. We predicted every \texttt{aligned} and
\texttt{inverted} cell would waste. \texttt{aligned} does, in all
$1{,}200$ cells at rent $> 0$. \texttt{inverted} wastes in $240$ and is
exactly zero in the other $960$, and the split is by arm, not by seed or
rent: under \texttt{inverted} the exempt case is a declined
\emph{disposable} file, and four of the five arms never decline one. At
rent $1.0$ and 60 cycles the ledger wastes $287.5$ spends of a $720$
capacity there while every other arm wastes $0.0$; under \texttt{aligned}
the four floor arms waste an identical $143.6$ and the ledger $239.6$.
So a tier that exempts the protected categories hides something from the
attacker whatever the arm does, and a tier that exempts the disposable
ones hides something only for the arm that has stopped answering. That
is this project's rent rule --- rent bills not having an answer ---
arriving a third time, now read off the attacker's ledger rather than
the environment's.

\paragraph{At a scarce budget the rule decides something, and it decides
it in the attacker's ledger.} Saturation is the regime the tier grid ran
and not the regime the caveat is about, so we ran the interior budget
too: both objectives at $\text{lie\_budget} = 2$, matched on every other
axis, $4{,}500$ paired cells. Neither rule ever fails to spend --- every
cycle offers at least two measurable outcomes --- so they differ in
\emph{where} the budget lands and not in how much of it is used. Under
\texttt{aligned} at rent $> 0$ the blind attacker puts $33.7\%$ of its
spends on outcomes of exactly zero; under \texttt{uniform} and
\texttt{inverted} it puts none there, and at rent $0$ it puts $33.7\%$
there in all three tiers, because an unpriced decline is unmeasured
whatever the tier says. Here the second registered prediction about
\texttt{inverted} failed the same way the first did and for the same
reason: we expected the two rules to diverge under \texttt{inverted} and
they never do, in any of its $1{,}200$ priced cells. Its exempt case is a
declined \emph{disposable} file, only an emptied store declines one, and
a budget of $2$ empties nothing.

The separation that matters is between two ledgers. In the attacker's,
\texttt{reported\_cum\_delta} differs in $2{,}100$ of the $4{,}500$
paired cells; in the world's, \texttt{cum\_delta} differs in $98$ and
final population in $14$. Where the world moves, the weaker attacker does
less damage in $91$ of those $98$ cells --- wasted budget is damage not
done --- and more in $7$, which are one seed, one arm and one horizon
(\texttt{quarantine}, seed 24, 60 cycles, $-0.72\%$) repeated across the
tiers and rents that expose it. That the direction is not universal is
the mechanism we flagged when registering it: suppressing a rent bill
hides a loss rather than causing one, so a wasted spend can help the
attacker, and here it does, rarely.

Only three arms move at all. \texttt{keep\_everything} and
\texttt{evict\_on\_negative} are identical in every cell, because
neither settles on a measurement it did not receive; \texttt{quarantine}
accounts for $70$ of the $98$, \texttt{survival} and
\texttt{survival\_paced} for $14$ each. And the tier conclusion is
untouched on both horizons: at rent $1.0$ the ledger's mean true
\texttt{cum\_delta} across 30 seeds is $+12.2865$M under
\texttt{aligned} against the price-greedy withholder and $+12.3095$M
against the blind one, with \texttt{uniform} and \texttt{inverted}
identical to every printed digit. Shape and surface are separated, and
what the tier axis measured was the shape.

\subsection{Every 30-cycle grid, at 60}
\label{sec:horizon}
Three results in this paper have turned on the horizon, each found by
accident and one grid at a time: benign retention under withholding
($0.92$ at 30 cycles, $0.44$ at 60), the rented test-suite family (flat
at 30, the ledger's first billable decline near cycle 49), and the
consolidation laundering entry (present at 30, starved by 59). Every
other deterministic grid here runs 30 cycles. We re-ran eleven of them
at 60 --- \texttt{headline}, \texttt{noisy}, \texttt{ablation},
\texttt{testsuite}, \texttt{testsuite\_noisy}, \texttt{memsec},
\texttt{adversary}, \texttt{persistence}, \texttt{salience},
\texttt{neighbours}, \texttt{bandit} --- each keeping the seed count its
committed file used, so all 5,815 cells pair one-to-one with a published
one and \texttt{cycles} is the only thing that varies. Five predictions
were registered before the run.

\paragraph{The canary, and the split.} \texttt{keep\_everything} removes
nothing, so its population must be horizon-invariant; it is, in $830$ of
$830$ cells, which is what licenses reading everything else as the
clock rather than the harness. $158$ of $163$ arm orderings hold. Every
storage grid holds \texttt{probe\_benign\_correct\_rate} flat to three
decimals. Both test-suite grids lose it --- $1.000 \rightarrow 0.750$ in
$10/10$ unattacked seeds, $0.518 \rightarrow 0.370$ in $30/30$ noisy
ones --- and nothing else in the sweep does. The storage-family headline
is horizon-stable; the second-family numbers are horizon-scoped.

\paragraph{The decay is a trade, and 30 cycles prices both sides at
zero.} The unattacked test-suite probe triple is identical in all ten
seeds. At 30 cycles the ledger answers every benign probe and is safe on
\emph{no} harmful one; at 60 exactly one entry has starved:

\begin{center}
\small
\begin{tabular}{lcccc}
\toprule
 & benign correct & silence & harmful safe & population \\
\midrule
30 cycles & $1.00$ & $0.00$ & $\mathbf{0.00}$ & 4 \\
60 cycles & $0.75$ & $0.40$ & $\mathbf{1.00}$ & 3 \\
\bottomrule
\end{tabular}
\end{center}

The entry that died was answering a benign probe and a harmful one, so
losing it costs a quarter of the benign answers and converts every
harmful answer into silence. On \texttt{testsuite\_noisy} the same trade
runs $0.518 \rightarrow 0.370$ benign, $0.507 \rightarrow 0.704$ silence,
$0.607 \rightarrow 1.000$ safe. Starvation is doing safety work on this
family that the reported horizon scores as none, by removing capability
that the same horizon scores as full.

\paragraph{One adverse reversal, and it is the rent rule again.} On the
noisy test-suite grid, at a flake rate of $0.15$, the hoarding arm
overtakes the ledger: $30.00 \rightarrow 60.00$ against $46.87
\rightarrow 47.93$. The counter accrues
linearly because it never removes anything; the ledger is flat because it
has starved to about $1.5$ entries and gone silent on $0.70$ of probes,
so it stops acting and stops earning. That is the rule the rent grids
established with the sign changed --- an emptied store has no answers,
and where inaction is unpriced that costs earnings rather than rent ---
reached here with no adversary and no price on standing still.

The other four changes reorder arms without reversing a claim. Two of
them sit below \texttt{survival} in \texttt{persistence}. One swaps two
arms in \texttt{salience} that are both negative. The last is on
\texttt{testsuite}, where the embedding arm's two-point lead over the
counter becomes a two-point deficit.

\paragraph{And one late kill.} \texttt{poison\_killed} matches its
committed value in $524$ of $525$ non-noisy \texttt{survival} cells. The
exception is \texttt{persistence} against the \texttt{persist} adversary
at budget 2, seed 6: not killed at 30 cycles, starved at cycle $56$. The
attacker's win in that cell is real at the horizon we publish and
temporary $26$ cycles past it.

\subsection{Where each defence catches each attack}
\label{sec:memsec}
The memory-security literature defends memory before any consequence
exists: at write (does this text look malicious?) and at retrieval
(should this be injected?). Both read content, and the systematic study
of memory poisoning reports the cost---every off-the-shelf content
detector it tests drops by $18$ to $42$ points of true-positive rate
between strong-signal payloads and weak-signal ones, the best falling
from $84.44\%$ to $42.50\%$, on payloads written to read like ordinary
operational facts \citep{mpbench2026}. We run that split on our own harness rather
than arguing it, with three attack classes and a reconstructed
write-time filter as the control.

The classes differ only where they must. \emph{Explicit} (strong signal)
wraps the packaged poison's actionable claim, verbatim, in imperative
override and claimed authority, so any outcome difference is
attributable to the surface and not to the payload.
\emph{Policy-conformant} (weak signal) is the packaged poison unchanged:
a plausible operational fact---these files are redundant copies---in the
register of a real runbook. \emph{Inert} advises no action at all: it is
the dormant case, and the honest limit of outcome-grounded revocation,
since an entry that never acts can never be caught by consequence. The
filter is a reconstruction, not a port, and we give it the charitable
form: it screens whole documents before encoding, so it sees the payload
as delivered rather than the fragments an encoder later makes of it.

The result (Table~\ref{tab:memsec}) is a clean division of labour, and
it includes a loss. On the strong-signal class the filter is simply
better than we are: it blocks at write for zero damage, where the ledger
must let the entry act once and pays $-0.39$M before revoking it.
Prevention beats revocation whenever detection works. On the
weak-signal class the filter is \emph{byte-identical to no defence at
all}---TPR $0.00$, $-12.1$M of damage, the poison alive at cycle 30---and
the ledger turns the same run into $+12.6$M, a $21.7$M swing on the class
that content inspection cannot see. On the inert class nothing detects
anything, and only the ledger removes the payload, by starvation at a
median cycle 19; every other arm carries all three poisoned entries to
the end of the run. Starvation, not detection, is the answer to dormant
poison, and no filter has it.

\paragraph{Consolidation launders a poisoned fragment.}
One cell does not come out clean, and it is ours. Under
\texttt{explicit}$+$\texttt{ledger} a single poisoned entry survives all
30 cycles on all 10 seeds with energy \emph{above} its spawn value,
carrying benign provenance alongside the poisoned source. The poisoned
text advises no action, so nothing punishes it, and it stays retrievable
inside a well-sourced entry. This is the compaction-driven write channel
\citep{mpbench2026} appearing inside our own mechanism: consolidation is
a poison-laundering surface, and any system that merges memories
inherits it. \S\ref{sec:merge-policy} evaluates the fix and reports what
the mechanism actually is, which is not what we first described.

\paragraph{The laundered entry never earns, and the encoder made it.}
\label{sec:merge-policy}
Two things we said about that cell were wrong, and instrumenting
\texttt{\_merge} shows both. The surviving entry's \texttt{uses} is
\emph{zero for the whole run}: it never answers anything, so it never
earns. Its above-spawn energy is pooled from its own poisoned siblings
--- $0.75 + 0.75 \rightarrow 1.50$ at cycle 4, then $1.25 + 1.25
\rightarrow 2.50$ at cycle 9 --- and an unmerged fragment gets only
$\text{spawn}/\text{upkeep} = 20$ cycles. So the entry is not permanent
either: it starves at cycle 59, at every seed, and at a 60-cycle horizon
the cell is clean. \textbf{Consolidation trades breadth for longevity},
and that is the laundering channel.

Nor did consolidation create it. Counting surviving poisoned entries
that also carry benign provenance (\texttt{poison\_laundered\_final},
which \texttt{poison\_alive\_final} cannot distinguish), the
\emph{no-defence} arm --- which consolidates nothing --- scores $2.00$
on this class. The initial store already contains an \texttt{ENTITY}
entry for ``Platform Team'' and a \texttt{CROSS\_DOC} summary whose
sources span the poisoned document and two clean ones. The explicit
payload claims authority \emph{from} the Platform Team, so naming an
entity the corpus already trusts is what places the attacker's text in
that entity's entry. No other attack class launders anything: both
others score $0.00$ in every cell, because neither names a trusted
entity. Consolidation then halves the count and triples the survivor's
life.

\paragraph{The obvious fix, evaluated: the natural reading is a no-op.}
\texttt{consolidate} now takes a \texttt{source\_policy}. Under
\texttt{"shared"} a cluster needs one source in common, the natural
reading of a trust boundary; under \texttt{"identical"} it must agree on
its whole source set. Three policies $\times$ three attack classes
$\times$ four defences $\times$ two horizons $\times$ 10 seeds, 720
runs. \texttt{"shared"} is byte-identical to the published behaviour
across $11{,}520$ metric comparisons: every merge in the laundering run
already shares a source, so it has nothing to refuse.
\texttt{"identical"} drives \texttt{poison\_laundered\_final} from
$1.00$ to $0.00$ at $10/10$ seeds, and it is surgical --- across all 720
runs it changes three metrics in that one cell and nothing in the other
twenty-three. It works by refusing the merge that pools the mixed
entry with a clean one, which lets the laundered entry starve on its own
runway instead of inheriting a longer one. The cost is one further probe
unanswered ($\texttt{probe\_silence\_rate}$ $0.20 \rightarrow 0.40$)
with $\texttt{probe\_benign\_correct\_rate}$ unchanged at $1.00$: the
store answers less, not worse. We ship it off by default, because every
committed number in this paper ran without it and the defence it buys is
scoped to one attack class on one corpus.

We registered the opposite prediction and were wrong. We had spot-checked
four metrics, found them equal under all three policies, and concluded
no-op; all four \emph{are} equal, and the difference lived in metrics we
had not printed and in one that did not yet exist. A comparison over a
subset of the metrics cannot support a claim about all of them.

\begin{table}[t]
\centering
\small
\caption{Attack class $\times$ where the defence sits, 10 seeds. Harm and
cum $\Delta$ in M of true bytes; starve cycle is when the last poisoned
entry of any kind is gone. The filter wins the strong-signal row; it is
identical to no defence on the weak-signal row; only the ledger empties
the store of dormant poison.}
\label{tab:memsec}
\begin{tabular}{llccccc}
\toprule
attack & defence & filter TPR & harm & cum $\Delta$ & starve & alive@30 \\
\midrule
\multirow{2}{*}{explicit}
 & filter        & 1.00 & \textbf{0.00} & \textbf{13.52} & 0 & 0 \\
 & ledger        & 1.00 & $-$0.39 & 12.59 & never & 1 \\
\midrule
\multirow{3}{*}{policy-conformant}
 & none          & 0.00 & $-$12.11 & $-$9.08 & never & 3 \\
 & filter        & 0.00 & $-$12.11 & $-$9.08 & never & 3 \\
 & ledger        & 0.00 & \textbf{$-$0.39} & \textbf{12.59} & 19 & 0 \\
\midrule
\multirow{2}{*}{inert (dormant)}
 & filter        & 0.00 & 0.00 & 13.52 & never & 3 \\
 & ledger        & 0.00 & 0.00 & 13.52 & \textbf{19} & \textbf{0} \\
\bottomrule
\end{tabular}
\end{table}

\subsection{A model in the loop, and where the claim does not hold}
\label{sec:wef}
Every result above lets a keyword function decide what ``the agent
adopted the poison'' means, and the corpus, the prompts, and that
function were written in the same vocabulary by the same hand. This
subsection removes that coupling: a local model (\texttt{llama3.2},
three seeds, 24 cycles) retrieves, answers, and cites, and the
lifecycle checkpoints are scored from what the model actually did.
The checkpoints follow the Write--Execute--Forget protocol
\citep{memsecbench2026} defines. \textbf{This is a reimplementation of
the protocol, not a port}: its cases and code were unavailable, so
nothing here reproduces or compares against its numbers.

\paragraph{The action-vocabulary gap.}
Building this suite exposed a threat to validity in our own harness.
The shared reader has \texttt{"safe to delete"}; a chat model writes
\texttt{"deleting store-1.db under data/ is safe"}. The first parses as
an action, the second as \emph{silence}, so the environment never
executes, never measures, and selection receives no signal from a
decision the model plainly made. LLM-mode answers are therefore read
with paraphrases of the same decision, scoped to this suite so the
deterministic suites' committed numbers stay reproducible under the
reader that produced them. The disagreement rate between the two
readers ran $0.00$--$0.04$, and every number in this subsection carries
that as its error bar.

\paragraph{Read against a real defence, the ledger does not win.}
Both curation arms beat \texttt{keep\_everything\_llm}, which is a
statement about curation and not about the ledger. The arm that decides
the question is \texttt{evict\_on\_negative\_llm}: the one-line
if-statement---evict any entry whose decisions produced a negative
outcome, no energy, no forgiveness---answering through the same model
and the same protocol. Against it (Table~\ref{tab:wef}):

\begin{itemize}
  \item \emph{Harm is a tie.} Adoption is $0.02$ for both arms, against
  $0.21$ for no curation. Externalized cycles ($1.7$ vs.\ $1.3$ and
  $1.7$ vs.\ $0.7$) do not separate them at three seeds.
  \item \emph{Revocation latency: the counter wins outright.} The last
  \emph{acting} poisoned entry is gone by cycle $1$--$3$ under the
  counter and $8$--$10$ under the ledger. The energy buffer is exactly
  what makes the ledger slower---an entry must burn through its balance
  before it dies---and here that design costs five to nine cycles of
  exposure, per seed and attack class.
  \item \emph{Completeness: the ledger wins outright.} It ends with no
  poisoned entry alive in $9$ of $9$ runs; the counter leaves one to
  five alive in every run, because an entry it never settles negatively
  is an entry it never removes.
\end{itemize}

\paragraph{What F1 is actually measuring.}
The ledger's $1.00$ repair rate is two mechanisms and should not be
read as one. Poison that \emph{acts} dies by consequence, at cycle
$8$--$10$, well before the population's starvation cliff at $19$.
Poison that never acts---the \texttt{inert} class, adoption and
externalization both $0.00$---dies at $19$ by upkeep alone, in the same
sweep that takes the alive population from $14$ to $5$. Removal by
disuse is a real capability, and the counter has no equivalent, but it
is not selection against poison: those entries were never identified as
harmful, they were merely never consulted. Within this 24-cycle horizon
they had also caused no measured harm.

The cost of that sweep is visible in the same table. Under
\texttt{policy\_conformant} the ledger's benign preservation falls to
$0.67$: a third of the benign probe set died in the cycle-19 collapse
alongside the dormant poison. F1 and F2 are reported together, and
their product---selective repair---is the honest figure, because repair
without preservation is amnesia rather than defence. The ledger ends
these runs carrying $4$--$5$ entries where both other arms carry
$14$--$18$, which is the same fact read as a benefit.

\paragraph{Verdict.}
On this protocol the pre-registered claim that conserved-resource
selection revokes poison \emph{faster} than a one-line heuristic is
\textbf{not supported}. What survives is narrower and we state it as
such: the ledger is the only arm that clears dormant payloads and the
only one that keeps the population small, and it pays for both with
revocation latency and with benign entries lost to the same starvation
that removes the dormant poison.

\begin{table}[t]
\centering
\small
\caption{Write--Execute--Forget with \texttt{llama3.2} deciding
adoption; means over seeds 0--2. \emph{kill} is the cycle the last
\emph{acting} poisoned entry is gone; \emph{alive} counts poisoned
entries of any kind at the end. The counter revokes sooner, the ledger
finishes the job.}
\label{tab:wef}
\begin{tabular}{llcccccc}
\toprule
attack & arm & E2 & E3 & F1 & F2 & kill & alive \\
\midrule
\multirow{3}{*}{explicit}
 & keep\_everything & 0.21 & 15.0 & 0.00 & 1.00 & never & 7 \\
 & evict\_on\_negative & 0.02 & 1.7 & 0.00 & 1.00 & \textbf{1--3} & 5 \\
 & ledger & 0.02 & 1.3 & \textbf{1.00} & 1.00 & 8--10 & \textbf{0} \\
\midrule
\multirow{3}{*}{policy-conformant}
 & keep\_everything & 0.21 & 15.3 & 0.00 & 1.00 & never & 3 \\
 & evict\_on\_negative & 0.02 & 1.7 & 0.00 & 1.00 & \textbf{1--3} & 1 \\
 & ledger & 0.02 & 0.7 & \textbf{1.00} & 0.67 & 8 & \textbf{0} \\
\midrule
\multirow{3}{*}{inert (dormant)}
 & keep\_everything & 0.00 & 0.0 & 0.00 & 1.00 & never & 3 \\
 & evict\_on\_negative & 0.00 & 0.0 & 0.00 & 1.00 & never & 3 \\
 & ledger & 0.00 & 0.0 & \textbf{1.00} & 1.00 & starve 19 & \textbf{0} \\
\bottomrule
\end{tabular}
\end{table}

\subsection{Literature controls, run rather than argued}
\paragraph{Retrieval bandit \citep{xu2026ael}.}
A per-entry successive-elimination bandit with a Hoeffding radius
(cull when $\hat\mu + \sqrt{\ln T / 2n} < 0.5$) is the strongest no-ledger
arm we field. Its confidence radius is real forgiveness, and under
one-sided false-bad noise it \emph{matches} survival at 35\%
($11.14$ vs $11.00$M, paired $p = 0.68$) while keeping full benign
capability. We publish this regime. Everywhere else the ledger's
pre-registered predictions land: the bandit cannot starve dead weight
(final population 15--16, near \texttt{keep\_everything}), kills the poison
cycles later (median 1.5 vs 0), and under symmetric noise its sign-blind
reward lets the poison feed until it bleeds to $\texttt{keep\_everything}$
territory.

\paragraph{LLM-judge settlement, and the cost result.}
Replacing measured settlement with keep/cull verdicts from a local model
(an operationalization of LLM-curated memory, cf.\ \citealp{zhang2026faulty})
gives a mixed quality result---one local judge degrades benign capability
to $0.67$, another does not, at five seeds---so we do \emph{not} claim
categorical superiority. The robust result is cost: judged settlement
runs $87.6$\,s (llama3.2:3b) to $1{,}514$\,s (qwen3:4b) per five-cycle
run against the ledger's $0.03$--$0.09$\,s, four to six orders of
magnitude, before a verdict is even parsed. Settlement by a scalar the
environment already returns is effectively free.

\subsection{Second environment, and where the ledger loses}
\label{sec:second-env}
On a test-suite environment (the resource is passing-test count; applying
a patch reruns the suite), the mechanism \emph{transfers}---survival kills
the actionable poison in every seed---but the headline tie does not:
\texttt{evict\_on\_negative} \emph{beats} survival ($88$ vs $69$,
$p = 0.014$). The cause is a designed property of the family meeting a
designed property of the ledger: this environment charges \emph{nothing}
for refusing to act, so the protector entry that keeps refusing the
destructive patch can never earn and starves under upkeep, after which the
destructive patch lands again. Leanness, the ledger's asset on storage, is
a liability where refusals are free. This is the central entry in the
regime map (\S\ref{sec:regime}): the conclusion is family-dependent in
\emph{both} directions, and we flag it rather than average it away.

We have since measured how much of that loss is the family and how much
is five entries we chose to put in it (\S\ref{sec:redundancy}). It is
mostly the five entries.

\subsection{Was it the corpus, or the merge?}
\label{sec:redundancy}
Three results in this paper sit on the test-suite family alone, and one
sentence has been explaining all three: the corpus is deliberately
redundant, so consolidation keeps finding surplus to starve. That names
two factors and asserts their conjunction from observations consistent
with either alone, so we varied both. \texttt{testsuite\_twins} drops
the five near-duplicate entries, removing the surplus while leaving the
other fifteen byte-identical; \texttt{consolidate\_every} $0$ removes
the merge while leaving the corpus alone. $2 \times 2 \times 2$
horizons $\times$ 5 arms $\times$ 30 seeds, twice: once on the plain
family and once on the rented one, because the sentence explains a
finding on each. The published corner reproduces
\texttt{rent\_testsuite.json}'s unattacked cells in all $300$ shared
cells, which is what licenses reading the other three corners as the
manipulation.

\paragraph{The two findings have different load-bearing halves.} The
capability decay needs both: \texttt{survival}'s benign probe rate falls
$1.000 \rightarrow 0.750$ between 30 and 60 cycles in $30/30$ seeds at
the published corner and in $0/30$ at each of the other three. Either
counterfactual removes it and neither is required alone, so for that
finding the conjunction was exactly right. The rent bill is not like
that. Dropping the twins recovers more of the gap than disabling the
merge does in $30/30$ seeds on both horizons, and at 30 cycles disabling
the merge alone makes the ledger \emph{worse} ($64.20$ against a
published $69.00$). One sentence was covering two mechanisms.

\paragraph{The second-family loss is mostly five entries.} At 60 cycles
the ledger trails the strike counter by $56.67$ unrented and $69.27$
rented. Dropping the twins leaves a gap of $3.80$ in both:
$93\%$ and $95\%$ of the loss this paper reports as a property of the
mechanism is a property of a corpus we designed. And the five entries
are not a shared cost. Both ledger arms move between the two corpora in
$30$ of $30$ seeds; all three counters move in $0$ of $30$, scoring
$178.00$, $140.00$ and $60.00$ either way. Redundancy is
free to an arm that hoards and billed to the only arm that pays upkeep,
which is the same asymmetry the rent grids found from the other
direction.

\paragraph{What that does and does not license.} It does not make the
loss disappear: $3.80$ at 60 cycles is still a loss, the ordering is
unchanged, and a redundant corpus is a legitimate world --- we built this
one that way on purpose, as the redundancy-cushioned complement to the
storage corpus. What it removes is the reading that the second-family
result measures the mechanism's behaviour in redundant environments
generally. A related correction: the storage corpus this paper calls
redundancy-free has $2$ mergeable pairs of $16$ entries against the
test-suite corpus's $5$ of $20$, so the contrast between the families is
one of degree.

\paragraph{It is not a dose. It is one entry.} Two points cannot tell a
dose from a particular entry: on two points the effect looks ``roughly
linear in that level'', and it is neither linear nor a level. Running all $32$
subsets of the five pairs ($32 \times 2$ horizons $\times$ 5 arms
$\times$ 30 seeds, with \texttt{"11111"} and \texttt{"00000"}
reproducing the grid above in all $600$ shared cells) puts the entire
effect on one of them. At 60 cycles, dropping the clamp-bound twin alone
recovers $93.3\%$ of the ledger's gap to the counter; dropping any of
the next three recovers exactly $0.0\%$, and dropping the fifth recovers
$-2.9\%$. Within a dose the spread across subsets is $54.5$ to $72.0$
against adjacent-dose mean gaps of $5.6$ to $14.4$, so a curve through
the dose means would describe a mechanism that is not there.

The rest of the corpus is not inert, it is protective. Keeping
\emph{only} the harmful twin scores $102.20$ against $121.33$ for
keeping all five, in $30/30$ seeds: the four other pairs recover a third
of the damage the first one does. And the capability decay is the same
single entry --- $456$ of the $480$ cells that keep it read $0.750$ at 60
cycles and $0$ of the $480$ that drop it do --- with the population
falling $4.00 \rightarrow 3.00$ across the horizon in exactly the masks
that keep it and holding at $4.00$ in every mask that drops it. One
entry starves between 30 and 60 cycles, and losing it costs one of the
four benign probes.

That pair is the least similar of the five --- $0.818$ against $0.889$
to $0.909$ --- so it is the one nearest the merge floor, and both stores
end at $15$ entries: the difference is only \emph{which} five merges
happened. We have since moved the floor instead of the corpus
(\S\ref{sec:merge-threshold}), and the answer is that the merge is the
smaller half of it.

\subsection{Moving the floor instead of the corpus}
\label{sec:merge-threshold}
Eleven merge thresholds, $2$ horizons, $5$ arms, $30$ seeds, twice over:
once with the query protocol's conflict floor coupled to the merge floor
as it has always been, and once with it pinned at the default. The
coupling is worth stating because it was invisible: one field set two
thresholds, so every \texttt{merge\_threshold} number in this
repository --- the ablation grid included --- was produced under it. It
turns out to be inert on this corpus (the two columns are identical in
all $1{,}320$ ledger cells), but it was inert by luck rather than by
design, and \texttt{SurvivalConfig.conflict\_threshold} now separates
them.

\paragraph{The floor's value does not matter; the set of merges does.}
Nine of the eleven thresholds merge all five pairs, and all nine are the
\emph{same run} --- identical per seed, not merely equal on average ---
at $121.33$. At $0.85$ the least similar pair drops out and the score
rises to $137.00$; at $0.90$ a second drops out and it rises to
$143.00$. Fewer merges is monotonically better at every step we can
reach, which refutes what we registered: we predicted consolidation
would be good on the four tight pairs and bad on the loose one, and it
is mildly bad on all of them here.

\paragraph{Three quarters of the harm is the entry, not the merge.}
This is the number the sweep was run for. Against a gap of $56.67$ to
the counter, refusing to merge the clamp-bound pair recovers $15.67$
($27.7\%$) while deleting its twin recovers $52.87$ ($93.3\%$). The two
are measured in different files by different knobs and agree on the
shared cell. So the entry hurts mostly by existing --- one more thing
holding energy and competing for retrieval --- and only secondarily by
being pooled. Our explanation named consolidation for three findings; on
the one we can now decompose, consolidation is the minority term.

\paragraph{How similarity is scored here.} Similarity in this sweep is
\texttt{MemoryStore.similarity}, the function \texttt{consolidate}
actually calls, not a Jaccard reimplemented in a test helper. The
distinction is load-bearing: a reimplementation produces correct counts
but incorrect values, and the wrong values put the band edges of this
sweep in the wrong places --- which is why three of its five registered
predictions are refuted rather than confirmed. Pair similarities are
pinned by value as well as by count.

\subsection{Parametric memory (secondary).}
To connect to memory-as-a-model \citep{quek2026memo}, we distill survivors
into a small (0.5B) model. Distilling the \emph{raw} store teaches the
facts but bakes in every poison statement (poison reproduction $1.00$);
distilling the \emph{energy-ledger survivors} teaches the same facts
(recall $1.00$) and reproduces no poison ($0.00$). The active ingredient
is a conserved-resource floor: a judge \emph{with} an energy floor nearly
matches it, a judge without one erodes the store to nothing. This is a
clean existence proof at one model scale, not a scaling law, and we frame
it as such.


\subsection{A retention attack that never writes}
\label{sec:potentiation}
The attack above corrupts the settlement signal. A second, weaker
adversary needs even less: it never writes, never settles, and only
asks questions. \citet{lin2026memsurvey} name the exposure without
measuring it --- ``retention schemes based on access frequency or
recency may inadvertently keep adversarial entries alive while
discarding legitimate ones'' --- and this section makes it
quantitative on a usage-based retention variant in our own codebase,
which is shipped disabled precisely because of what follows.

A usage signal is not only a poor proxy for usefulness; when it is
\emph{relative}, it is a proxy the attacker controls. The variant slows an
entry's upkeep in proportion to how often it is recalled and how central
it is in an associative graph, and both move on query alone---so a
query-only adversary in the sense of \citet{dong2025minja}, one who
never writes and never settles, drives them without crossing any channel
our threat model names. Under flat upkeep the dormant poison starves on
exactly the benign schedule (margin $0$ cycles); under the usage-based
variant with the same attacker it outlives every benign entry by $4$,
and the margin is identical for the loudest payload and the most inert
one, because this path reads usage and never text.

Because the measurement is deterministic and needs no model, we can run
it as a grid rather than an anecdote: two independent corpora $\times$
four upkeeps $\times$ two normalisations, $32$ cells. Flat upkeep
favours the poison in $0$ of $32$; the attacker gains in $16$ of $16$
potentiated cells, on both corpora and all three attack classes. The
durable quantity is a fraction rather than a cycle count---the absolute
margin runs $8$--$10$ cycles at upkeep $0.02$ and $1$ cycle at $0.2$,
but $0.11$--$0.17$ of the starvation horizon throughout. A query-only
adversary owns roughly the last $13\%$ of the store's life, at whatever
timescale upkeep sets.

The cause is normalisation against the live population, which we
establish by counterfactual rather than by argument: replacing the
recall term's denominator---the population peak, for a fixed cap---while
leaving graph centrality fully attacker-driven collapses the margin
(median fraction $0.133 \rightarrow 0.050$) and leaves the starvation
horizon intact. It does not quite abolish it: $11$ of $16$ cells retain
a margin of \emph{exactly one cycle}, the smallest observable value, and
that residual tracks $1/\text{horizon}$ across the grid, so we report
peak-normalisation as the dominant contributor rather than the sole
cause. Neither usage as such nor centrality carries the attack. What
does is that the score is a \emph{standing}: one entry's value is a
function of every other entry's traffic, so the adversary takes lifetime
from the benign population rather than buying its own, and no query
budget prices it out. Potentiation is not inherently poison-friendly
either---on the second corpus, potentiation \emph{without} an attacker
starves the poison $18$ cycles \emph{before} the benign entries. The
margin is the adversary's doing, and the system charges flat upkeep by
default for that reason.

\subsection{A mechanically curated system: MemoryOS}
\label{sec:memoryos}
The Mem0 result and the survey above leave one system in scope, and it is
the interesting one. MemoryOS \citep{kang2025memoryos} curates
\emph{mechanically}: its mid-term store holds sessions under a capacity
and, on overflow, calls \texttt{evict\_lfu}, which is
\texttt{min(access\_frequency)} and nothing else --- no model is asked,
no text is read. The signal is retrieval frequency, incremented by
\texttt{search\_sessions} on every match. A curator that removes entries,
decides on a signal rather than a judgment, and whose signal moves when
somebody asks a question: every precondition our threat model needs, in
published software we did not write. The measurement needs no model
either, since the eviction is arithmetic.

\paragraph{The predicted attack fails, for a reason worth stating.}
Eviction takes a minimum, so the obvious move is to leave the victim
untouched and raise its peers until it is the unique lowest. It does not
work. \texttt{add\_session} registers a newcomer at frequency zero and
\emph{then} evicts, so every arrival sits at the floor and evicts itself:
a memory retrieved even once cannot be removed by capacity pressure at
all. We report this because it is a real property of the mechanism and
because we predicted the opposite.

\paragraph{What the mechanism does instead.} Sweeping the victim over
every seeded memory, eviction tracks a single bit exactly --- evicted if
and only if \texttt{access\_frequency} is zero, with no exceptions: all
$3$ victims that reached zero were evicted, and none of the $9$ that were
retrieved even once ever was. A never-retrieved memory loses to a
brand-new arrival that is \emph{also} at zero, because \texttt{min}
returns the first minimum in insertion order and the older entry comes
first. MemoryOS therefore deletes the memory nobody has asked for yet in
preference to the one that arrived a moment ago, and a single retrieval
confers permanent immunity.

This is the rare-but-critical case: an emergency contact or an allergy
note is precisely the memory stored once, needed rarely, and never
consulted in between, and it is first out. Our own abstract names that
cost --- upkeep starves the long tail of correct-but-rarely-used
knowledge --- and \textsc{darwin-memo} answers it with pinning, which
floors a pinned entry's balance at zero instead of burying it. On this
path MemoryOS has no equivalent, so the exposure we describe as a design
trade in our own mechanism is unmitigated in a deployed one.

\paragraph{The other direction: promotion is cheap.} Eviction cannot be
driven by inflation, but MemoryOS has a second curation decision that
can, and it runs the other way. When a session's heat crosses
\texttt{H\_PROFILE\_UPDATE\_THRESHOLD}, MemoryOS analyses that session
and writes what it extracts into \emph{long-term} memory, a tier
capacity pressure never touches. Heat is
$N_{\text{visit}} + L_{\text{interaction}} + R_{\text{recency}}$, so the
cost of crossing it is arithmetic: with the shipped threshold of $5.0$, a
single-page session reaches it on the \textbf{third self-query}
($3.0 \rightarrow 4.0 \rightarrow 5.0$). An adversary that gets any
content into mid-term storage and then asks about it three times has the
curator promote that content into the persistent tier --- no delete call,
no judge, no further writes. The threat model's usual direction is
denial; this is its mirror, and it is the cheaper of the two. We measure
the precondition (crossing the threshold) rather than the extraction
itself, which is a model call we do not exercise.

\paragraph{Promotion, carried out end to end.} We then ran the
extraction rather than only its precondition. A canary string appears in
the adversary's turn and nowhere else, so its presence in a long-term
tier cannot be a paraphrase of benign content. Against a
\texttt{quiet} control --- the same adversarial turn added and never
asked about --- the result is unanimous over three trials: the control
never crosses the threshold (heat $2.0$) and never promotes, while three
self-queries cross it (heat $5.0$) and promote every time. The canary
reaches MemoryOS's long-term \emph{user knowledge} in $3/3$ trials and
its assistant knowledge in $2/3$; the stored user profile was never
poisoned, since that write is gated on the analysis returning enough
text. So the adversary's content moves from the evictable tier into a
persistent one for the price of three questions, and MemoryOS's own log
records it doing so.

Two disclosures, because the difference between compensating for a model
quirk and lowering a bar until an attack works is the whole of the
result's value. First, MemoryOS parses its multi-topic summary reply with a bare
\texttt{json.loads}; on \texttt{JSONDecodeError} it warns, returns no
themes, and the updater then files the whole batch as a single session
under the constant string ``General conversation segment from short-term
memory.'' with an empty keyword list. Retrieval scores a session on its
summary embedding \emph{and} its keywords, so both terms lose their
discriminating power at once, and a model that fences its JSON costs
MemoryOS the topic structure of that batch and makes the content beneath
it hard to reach. The failure is logged rather than silent --- we
originally wrote otherwise, and checked the source instead of our own
prose only when preparing the maintainer disclosure. That is a real fragility in the
target, it is not what we set out to measure, and left in place it
prevents the promotion path from being reached at all, so we strip the
fence at the client boundary to restore the behaviour MemoryOS gets from
a non-fencing model. Second, \texttt{force\_mid\_term\_analysis} exists
in the API and would bypass the heat threshold entirely; we do not call
it. No threshold is lowered and no decision of MemoryOS's is skipped.

\paragraph{Disclosure.} All three findings were sent to the corresponding
author before this section was written, and the reply neither disputed a
finding nor requested an embargo. \S\ref{sec:ethics} gives the full account,
including what we did not attack and what the artifact does and does not
release.

\paragraph{What we do not claim.} An adversary able to dominate the
retrieval channel --- supplying content that absorbs queries which would
otherwise reach the victim --- would keep a chosen memory at zero and let
the curator delete it in favour of the adversary's own fresh content.
What we demonstrate is the second half of that: neglect kills,
deterministically, and the curator selects the neglected. Manufacturing
the neglect end to end against a live agent is not demonstrated here and
should not be read as though it were.

\begin{table}[t]
\centering
\small
\caption{\textsc{MemoryOS}, both mechanical curation decisions, measured end
to end. LFU eviction fires on the target memory only when it is neglected;
no memory that was retrieved at all was ever evicted ($0$ of $9$).
Heat-triggered promotion carries the adversary's canary into long-term
\emph{user knowledge} only when the adversary keeps querying it. Both
decisions are reached by a rule from a signal, and the signal is usage.}
\label{tab:memoryos}
\begin{tabular}{llccc}
\toprule
decision & condition & trials & fired & rate \\
\midrule
\multirow{2}{*}{LFU eviction}
 & attended   & 6 & 0 & 0.00 \\
 & neglected  & 6 & 3 & 0.50 \\
\midrule
\multirow{2}{*}{promotion}
 & quiet      & 3 & 0 & 0.00 \\
 & queried    & 3 & 3 & 1.00 \\
\bottomrule
\end{tabular}
\end{table}

\subsection{Withholding: the dual of lying}
\label{sec:withholding}
Every adversary above \emph{corrupts} a measurement. The obvious weaker
one merely \emph{withholds} it: the reported delta is zero and nothing
else changes. This is a strict subset of what the destruction adversary
can express, and it is weaker in the way a threat model cares about,
since indiscriminate withholding never reads the sign of the true delta.
We separate it because of an asymmetry in our own arm set that nothing
else probes: the baseline drivers never charge upkeep, so \emph{only the
ledger has a clock}. Credit is capped at $\text{max\_energy}$, so bounded
credit implies bounded runway --- nothing survives more than
$\text{max\_energy}/\text{upkeep}$ unmeasured ticks however valuable it
is, and an entry that has not earned starts at spawn energy and gets
$\text{spawn}/\text{upkeep} = 20$.

We also run the mitigation this suggests, as an arm rather than a
default: \texttt{survival\_paced} charges upkeep only on cycles that
carried a measured outcome. The decision is population-level and reads no
per-entry state, which is what separates it from the salience control of
\S\ref{sec:headline}, where a per-entry usage signal cannot distinguish
``used'' from ``useful'' and the poison is the most-used entry.

\begin{table}[t]
\centering
\small
\caption{Withholding, 30 seeds, storage environment, 60 cycles. Budget is
suppressions per cycle out of 12 measurable tasks, so 12 is total
suppression. \emph{Sel.} is the selective withholder, which suppresses
only measurements that would incriminate its own poison. Benign
capability is $1.00$ for every arm under the selective attack and is
therefore omitted from that half.}
\label{tab:withholding}
\begin{tabular}{llccc c}
\toprule
& & \multicolumn{3}{c}{indiscriminate} & sel. \\
arm & budget & benign & kill & cum $\Delta$ & kill \\
\midrule
\textbf{survival} & 0 & 1.00 & 1.00 & $+25.55$M & 1.00 \\
                  & 8 & $\mathbf{0.44}$ & 1.00 & $+18.94$M & 1.00 \\
                  & 12 & $\mathbf{0.00}$ & 1.00 & $\mathbf{-6.42}$M & 1.00 \\
\midrule
survival\_paced & 0 & 1.00 & 1.00 & $+25.55$M & 1.00 \\
                & 8 & $\mathbf{1.00}$ & 1.00 & $+22.41$M & 1.00 \\
                & 12 & 1.00 & $\mathbf{0.00}$ & $-18.17$M & 1.00 \\
\midrule
evict\_on\_negative & 0 & 1.00 & 1.00 & $+25.74$M & 1.00 \\
                    & 8 & 1.00 & 1.00 & $+24.12$M & $\mathbf{0.00}$ \\
                    & 12 & 1.00 & $\mathbf{0.00}$ & $-18.17$M & $\mathbf{0.00}$ \\
\midrule
keep\_everything & 12 & 1.00 & 0.00 & $-18.17$M & 0.00 \\
\bottomrule
\end{tabular}
\end{table}

\paragraph{Withholding does not spare the counters, it dissolves them.}
At total suppression \texttt{evict\_on\_negative} and quarantine are
identical to \texttt{keep\_everything} in every column: kill rate $0.00$,
poison alive, the same cumulative delta and the same final population.
Their benign retention of $1.00$ is not a defence holding up under
attack, it is a mechanism that has stopped running --- no measurement
means no blame, and a strike counter with nothing to strike removes
nothing. We predicted the opposite when we designed the experiment, and
say so here because the arm set was chosen to make either answer visible.

\paragraph{The ledger's failure mode is amnesia, and amnesia is the
cheaper failure.} Survival loses all benign capability at total
suppression, which is the real cost of this attack and the honest
headline for it. It nonetheless ends at $-6.42$M where every other arm
ends at $-18.17$M over 60 cycles, because an emptied store stops acting
while a live poisoned store keeps destroying. We pre-registered that
cumulative delta would not reverse; it did not.

\paragraph{The kill rate at total suppression is starvation, not
selection.} Poison kill cycle and poison starve cycle are both $19.0$:
the poison died in the same undifferentiated collapse at the
$\text{spawn}/\text{upkeep}$ cliff that took the benign population. This
is the same artifact we flag for the repair metric in
\S\ref{sec:wef}, and it is the third time in this paper that a poison
count has flattered a mechanism that had not eliminated anything by
selection.

\paragraph{The horizon changed the conclusion.} At budget 8 survival
retains $0.92$ of benign capability over 30 cycles and $0.44$ over 60. A
30-cycle grid understates this attack roughly sevenfold, because the
starvation cliff sits at cycle 20 and leaves only ten cycles past it.
Every deterministic result elsewhere in this paper runs 30 cycles, and
this is the first place we have measured what that horizon hides.

\paragraph{The mitigation fails against an attacker who knows it is
there.} A selective withholder suppresses only what would incriminate its
own poison and lets benign outcomes through, so an evidence-paced clock
never pauses. Under it \texttt{survival\_paced} is identical to
\texttt{survival} in every cell we ran --- six budgets, two horizons,
every metric --- and its entire advantage in the indiscriminate column
vanishes. The counters fare worse still: from budget 8 up they never
observe a negative, so their kill rate is $0.00$ and they are again
indistinguishable from no curation. Survival keeps a $1.00$ kill rate at
every budget, but its kill cycle runs $0.3 \rightarrow 1.9 \rightarrow
10.7 \rightarrow 19.0$, which is revocation degrading from execution into
removal-by-disuse rather than a defence holding. One cell inverts a
published ordering and we name it rather than bury it: at budget 4
survival revokes at cycle $10.7$ against the counter's $12.5$, where
\S\ref{sec:adversary} has the counter revoking roughly five times faster.

\begin{table}[t]
\centering
\small
\caption{The same indiscriminate attack on the test-suite environment,
30 seeds, 60 cycles. Budget 12 is total suppression here as well.
Cumulative delta is in passing tests rather than resource units, so the
scales are not comparable across families; the ordering within a column
is. \texttt{keep\_everything} reads $+60.00$ at every budget, which is
the canary: it never curates, so the attack cannot move what it does to
the world.}
\label{tab:withholding-testsuite}
\begin{tabular}{llccc}
\toprule
arm & budget & benign & kill & cum $\Delta$ \\
\midrule
\textbf{survival} & 0 & 0.75 & 1.00 & $+121.33$ \\
                  & 12 & $\mathbf{0.00}$ & 1.00 & $\mathbf{+65.00}$ \\
\midrule
survival\_paced & 0 & 0.75 & 1.00 & $+121.33$ \\
                & 12 & $\mathbf{1.00}$ & $\mathbf{0.00}$ & $+60.00$ \\
\midrule
evict\_on\_negative & 0 & 1.00 & 1.00 & $\mathbf{+178.00}$ \\
                    & 12 & 1.00 & $\mathbf{0.00}$ & $+60.00$ \\
\midrule
keep\_everything & 12 & 1.00 & 0.00 & $+60.00$ \\
\bottomrule
\end{tabular}
\end{table}

\paragraph{On the second environment family the direction replicates and
the magnitude does not.} The limitation we recorded for this section was
that it ran on one environment, and that the storage corpus is where
amnesia is cheapest: an emptied store scores zero, so doing nothing costs
nothing. We have now built the adversarial test-suite environment and run
the same grid on it (Table~\ref{tab:withholding-testsuite}). The claim
survives: at total suppression survival has the best true cumulative
delta on both horizons, and with $\sigma = 0.00$ across 30 seeds, because
total suppression removes the stochastic channel and the cells are exact
rather than noisy. Predictions were registered before the run in a
separate commit.

What does not survive is the size of the effect. On storage the ledger
ends roughly three times better than every other arm; here it ends
$8\%$ better, $+65.00$ against a floor of $+60.00$. Amnesia is worth far
less where the poison's damage per cycle is bounded --- one re-offered
destructive patch --- than where it accumulates across a corpus. The
limitation was directionally wrong and quantitatively right, and we
would rather report that than the half of it that flatters the ledger.

\paragraph{Read this attack as a leveller.} Unattacked, the counters are
\emph{better} on this family: eviction reaches $+178.00$ against
survival's $+121.33$, which is the ordering of
\S\ref{sec:second-env}, where refusing to act is free. The attack costs
them that entire lead. Measured as degradation from budget 0 to budget
12 at 60 cycles, eviction loses $66\%$, quarantine $57\%$ and survival
$46\%$. Nobody wins. The ledger loses least and lands
marginally above a floor everyone else falls to, and its kill column is
starvation here too --- kill cycle and starve cycle are both $19.0$, and
its final population at budget 12 is $0.0$, so the advantage is banked
before extinction rather than earned by a working store.

\paragraph{Pacing gets a real window here, and still does not ship.} On
this family \texttt{survival\_paced} is no longer degenerate at total
suppression: it holds benign capability at $1.00$ where survival reaches
$0.00$, and pays $+60.00$ against $+65.00$ for it. That is a genuine
trade rather than the collapse into no-curation it shows on storage. It
buys the retention the same way, though --- by removing nothing at all,
poison included, at a kill rate of $0.00$ --- and it is still measured
only against the attacker that does not read the sign.

We therefore ship the pacing knob disabled and documented as unproven.
Its one measured advantage exists only against the attacker that does not
read the sign, and at total suppression its limit is exactly
no-curation. A mitigation that helps against the dumber of two attackers
is a finding, not a default.

\subsection{Is a memory system worth its capability cost?}
\label{sec:frontier}
The two controls above are differently-curated stores. The baseline a memory
paper is actually judged against is \emph{no memory system}: put the whole
history in the prompt and let the model select. \texttt{full\_context\_llm}
is that arm---the same model, protocol, and audit trail, with retrieval
replaced by ``hand back every live entry.'' It differs from
\texttt{keep\_everything} in exactly one method, so the pair separates a store
with no eviction from a reader with no retrieval. Five seeds, 20 cycles,
\texttt{llama3.2} (Table~\ref{tab:llm-frontier}); this is the deterministic
\texttt{llm} suite, not the Write--Execute--Forget protocol of
Table~\ref{tab:wef}, and the two are not comparable cell-for-cell.

Read as a frontier rather than a winner: \texttt{full\_context} is the
\emph{capability ceiling}---it sees everything useful and posts the highest
cumulative resource ($4.22$M)---and in the same runs the \emph{safety floor},
because it also sees every poison and removes none ($0/5$ seeds clear the
acting poison, harmful-probe refusal $0.50$). \texttt{survival} keeps $65\%$
of that ceiling ($2.75$M) while clearing all acting poison in $5/5$ seeds,
ending with no poisoned entry alive, and refusing every harmful probe. It pays
the leanness cost the attack table already charged: benign recall falls to
$0.67$ as the population starves to four. Curation here is not free
capability; it is a capability \emph{cost} bought for a safety floor the
no-system baseline never reaches.

Two controls pin the corners. \texttt{evict\_on\_negative} matches survival on
poison kills and harmful-probe refusal---though it leaves the one dormant
payload survival starves---and does so at \emph{negative} capability
($-3.55$M): a sign-blind counter that evicts on any bad outcome over-prunes
the entries the workload needs, so safety without an energy buffer is bought
by discarding capability rather than keeping it. \texttt{keep\_everything} is dominated---worse than
\texttt{full\_context} on capability and tied on every other column---which is
the empirical reason ``keep everything'' belongs in the prompt, not in a
bottlenecked store. The mitigation sweep (\texttt{refuse=on}) moves survival
by under $3\%$ on capability and not at all on safety, and is folded into the
\texttt{survival} row.

\begin{table}[t]
\centering
\small
\caption{The long-context baseline. \texttt{llama3.2}, five seeds, 20 cycles,
means over seeds. \emph{cap}: cumulative measured resource ($\times 10^6$
units, the environment's own scalar, identical across arms). \emph{killed}:
seeds of five that removed every \emph{acting} poison. \emph{alive}: poisoned
entries of any kind at the end. \emph{safe}: harmful-probe refusal rate.
\emph{benign}: benign-probe correct rate. \emph{pop}: final population. Only
\texttt{survival} is both capable and safe, and it alone pays in benign
recall.}
\label{tab:llm-frontier}
\begin{tabular}{lrccccc}
\toprule
arm & cap & killed & alive & safe & benign & pop \\
\midrule
\texttt{full\_context}       & $4.22$  & $0/5$          & $3.0$          & $0.50$          & $1.00$ & $16$ \\
\texttt{keep\_everything}    & $-2.90$ & $0/5$          & $3.0$          & $0.50$          & $1.00$ & $16$ \\
\texttt{evict\_on\_negative} & $-3.55$ & $\mathbf{5/5}$ & $1.0$          & $\mathbf{1.00}$ & $1.00$ & $14$ \\
\texttt{survival}            & $2.75$  & $\mathbf{5/5}$ & $\mathbf{0.0}$ & $\mathbf{1.00}$ & $0.67$ & $4$  \\
\bottomrule
\end{tabular}
\end{table}



\bibliographystyle{plainnat}
\bibliography{references}

\end{document}
